%%%%%%%% ICML 2025 EXAMPLE LATEX SUBMISSION FILE %%%%%%%%%%%%%%%%%

\documentclass{article}

% Recommended, but optional, packages for figures and better typesetting:
\usepackage{microtype}
\usepackage{graphicx}
\usepackage{subfigure}
\usepackage{booktabs} % for professional tables
% The ICML style loads xcolor itself. Pass the table option before loading the
% style to avoid an option clash while retaining colored table support.
\PassOptionsToPackage{table}{xcolor}

% hyperref makes hyperlinks in the resulting PDF.
% If your build breaks (sometimes temporarily if a hyperlink spans a page)
% please comment out the following usepackage line and replace
% \usepackage{icml2025} with \usepackage[nohyperref]{icml2025} above.
\usepackage{hyperref}


% Attempt to make hyperref and algorithmic work together better:
\newcommand{\theHalgorithm}{\arabic{algorithm}}

% Use the following line for the initial blind version submitted for review:
% \usepackage{icml2025}

% If accepted, instead use the following line for the camera-ready submission:
\usepackage[accepted]{icml2025}
% \usepackage{icml2025}

% For theorems and such
\usepackage{amsmath}
\usepackage{amssymb}
\usepackage{mathtools}
\usepackage{amsthm}
\usepackage{algorithm}
\usepackage{algorithmic}
\usepackage{wrapfig}
\usepackage{tabularx} 
\newcolumntype{P}[1]{>{\centering\arraybackslash}p{#1}}

% if you use cleveref..
\usepackage[capitalize,noabbrev]{cleveref}

%%%%%%%%%%%%%%%%%%%%%%%%%%%%%%%%
% THEOREMS
%%%%%%%%%%%%%%%%%%%%%%%%%%%%%%%%
\theoremstyle{plain}
\newtheorem{theorem}{Theorem}[section]
\newtheorem{proposition}[theorem]{Proposition}
\newtheorem{lemma}[theorem]{Lemma}
\newtheorem{corollary}[theorem]{Corollary}
\theoremstyle{definition}
\newtheorem{definition}[theorem]{Definition}
\newtheorem{assumption}[theorem]{Assumption}
\theoremstyle{remark}
\newtheorem{remark}[theorem]{Remark}

% Todonotes is useful during development; simply uncomment the next line
%    and comment out the line below the next line to turn off comments
%\usepackage[disable,textsize=tiny]{todonotes}
\usepackage[textsize=tiny]{todonotes}

\newcommand{\jhy}[1]{\textcolor{red}{#1}}
\newcommand{\lhl}[1]{\textcolor{blue}{#1}}
\newcommand{\shortname}{CateKV }
\newcommand{\jc}[1]{\textcolor{teal}{#1}}

\definecolor{comment_color_2}{RGB}{64,128,128}
\newcommand{\LineComment}[1]{\vspace*{0.5em}\small\textcolor{comment_color_2}{\textit{\# #1}}}

% The \icmltitle you define below is probably too long as a header.
% Therefore, a short form for the running title is supplied here:
\icmltitlerunning{CateKV: On Sequential Consistency for Long-Context LLM Inference Acceleration}

\begin{document}

\twocolumn[
\icmltitle{CateKV: On Sequential Consistency for Long-Context LLM Inference Acceleration}


\icmlsetsymbol{equal}{*}

\begin{icmlauthorlist}
\icmlauthor{Haoyun Jiang}{sjtu,ali}
\icmlauthor{Haolin Li}{fdu}
\icmlauthor{Jianwei Zhang}{ali}
\icmlauthor{Fei Huang}{ali}
\icmlauthor{Qiang Hu}{sjtu}
\icmlauthor{Minmin Sun}{ali}
\icmlauthor{Shuai Xiao}{ali}
\icmlauthor{Yong Li}{ali}
\icmlauthor{Junyang Lin}{ali}
\icmlauthor{Jiangchao Yao}{sjtu}
\end{icmlauthorlist}

\icmlaffiliation{sjtu}{CMIC, Shanghai Jiao Tong University}
\icmlaffiliation{ali}{Alibaba Group}
\icmlaffiliation{fdu}{Fudan University}

\icmlcorrespondingauthor{Jiangchao Yao}{sunarker@sjtu.edu.cn}
\icmlcorrespondingauthor{Shuai Xiao}{shuai.xsh@gmail.com}

\icmlkeywords{Machine Learning, ICML}

\vskip 0.3in
]

% this must go after the closing bracket ] following \twocolumn[ ...

% This command actually creates the footnote in the first column
% listing the affiliations and the copyright notice.
% The command takes one argument, which is text to display at the start of the footnote.
% The \icmlEqualContribution command is standard text for equal contribution.
% Remove it (just {}) if you do not need this facility.

\printAffiliationsAndNotice{}  % leave blank if no need to mention equal contribution
% \printAffiliationsAndNotice{\icmlEqualContribution} % otherwise use the standard text.

\begin{abstract}
Large language models (LLMs) have demonstrated strong capabilities in handling long-context tasks, but processing such long contexts remains challenging due to the substantial memory requirements and inference latency. In this work, we discover that certain attention heads exhibit sequential consistency in their attention patterns, which can be persistently identified using a coefficient-of-variation-based algorithm. Inspired by this observation, we propose CateKV, a hybrid KV cache method that retains only critical token information for consistent heads, thereby reducing KV cache size and computational overhead, while preserving the majority of KV pairs in adaptive heads to ensure high accuracy. We show the unique characteristics of our algorithm and its extension with existing acceleration methods. Comprehensive evaluations on long-context benchmarks show that, while maintaining accuracy comparable to full attention, CateKV reduces memory usage by up to $2.72\times$ and accelerates decoding by $2.18\times$ in single-sample inputs, and boosts throughput by $3.96\times$ in batch scenarios. The code is available at \url{https://github.com/haoyun-jiang/CateKV}.
% Long-context large language models (LLMs) are essential for applications requiring extensive contextual understanding. However, processing such long contexts significantly hampers inference speed and increases memory usage due to the need to store and retrieve large key-value (KV) caches during self-attention. In this work, we discover that certain attention heads exhibit sequential consistency in their attention patterns and classify them into consistent and adaptive categories using a coefficient of variation–based algorithm. Our hybrid KV cache strategy retains only critical token information for consistent heads, thereby reducing KV cache size and computational overhead, while adaptive heads preserve the majority of KV pairs to ensure high accuracy during decoding. Comprehensive evaluations on benchmarks like RULER, Longbench, and NIAH demonstrate that our method achieves up to $n\times$ reduction in memory usage and $m\times$ acceleration in decoding speed on a single NVIDIA A100 GPU, while maintaining performance comparable to full attention mechanisms. Additionally, when integrated with other optimization techniques, our framework realizes an $a\times$ speedup, showcasing its effectiveness and scalability for deploying LLMs in long-context scenarios.
\end{abstract}

\section{Introduction}
% The rapid development of large language models (LLMs) has significantly impacted human life. With the growing demand for multi-turn dialogues and long-context inputs, many contemporary models now support context windows of 128K tokens or more~\cite{gpt,llama,qwen,chatglm,gemini,phi}. This capability enables effective performance in tasks involving lengthy texts, such as long-document question answering~\cite{peek,multiqa}, extensive information retrieval~\cite{infinitebench,informationretrieval}, and repository-level code understanding~\cite{codeplan,swebench}. However, the autoregressive nature of transformer-based models leads to increased computational costs, memory access overhead, and memory consumption as context lengths grow, due to the need to store and retrieve all key-value (KV) caches. For instance, using the Llama-3-8B-Instruct-262k~\cite{llama3-8b-262k} model with FlashAttention, extending the context from 1K to 1M tokens increases end-to-end inference latency by more than 3000 times~\cite{minference}. Moreover, when processing an input of approximately 500K tokens with the Llama-3-8B model, memory consumption exceeds 80GB, surpassing the capacity of a single NVIDIA A100-80GB GPU. Consequently, accelerating long-context LLMs is both essential and imperative.

% With the powerful capabilities of large language models (LLMs), there is an increasing demand for long-context inputs. Many contemporary models now support context windows of 128K tokens or more~\cite{gpt,llama,qwen,chatglm,gemini,phi}, enabling effective performance in tasks like long-document question answering~\cite{peek,multiqa}, information retrieval~\cite{infinitebench,informationretrieval}, and repository-level code understanding~\cite{codeplan,swebench}. However, the autoregressive nature of transformer-based models leads to increased computational costs, memory access overhead, and consumption as context lengths grow, due to the need to store and retrieve all key-value (KV) caches. For example, using the Llama-3-8B-Instruct-262k~\cite{llama3-8b-262k} model with FlashAttention, extending the context from 1K to 1M tokens increases inference latency by over 3000 times~\cite{minference}, while processing 550K tokens with the Llama-3-8B model exceeds 80GB in memory usage, surpassing the capacity of a single NVIDIA A100-80GB GPU. Consequently, accelerating long-context LLMs is both essential and imperative.


With rapid development of large language models (LLMs), many generalist models support context windows of 128K tokens or more~\cite{gpt,llama,qwen,chatglm,gemini,phi}, enabling to effectively perform tasks like long-document question answering~\cite{peek,multiqa}, information retrieval~\cite{infinitebench,informationretrieval}, and repository-level code understanding~\cite{codeplan,swebench}. However, as context lengths grow, the autoregressive nature of the mainstream LLM paradigm often leads to increased computational costs, memory consumption, and thus the runtime, since we have to store and retrieve all key-value (KV) caches. For example, using the Llama-3-8B-Instruct-262k~\cite{llama3-8b-262k} model with FlashAttention, extending the context from 1K to 1M tokens increases inference latency by over 3000 times~\cite{minference}.
Therefore, accelerating LLM inference in long contexts is both essential and imperative.

\begin{figure}[t]
\begin{center}
    \includegraphics[width=0.48\textwidth]{Figures/intro.pdf}
    % \vspace{-0.1cm}
    \caption{Comparison of KV eviction, KV retrieval, and our method. CateKV employs a hybrid KV cache method to reduce memory usage while maintaining high accuracy.
    % \lhl{put the explanation to the main body}
    % The KV Eviction Algorithm suffers from accuracy loss due to irreversible eviction, while KV Retrieval incurs high memory usage by retaining all KV pairs. In contrast, our algorithm balances accuracy and memory by retaining critical KV pairs in consistent heads and most KV pairs in adaptive heads for selective retrieval during decoding. 
    }
    \vspace{-0.4cm}
    \label{fig:intro1}
\end{center}
\end{figure}

Existing methods for inference acceleration of LLMs under long contexts can be categorized into two types: KV cache eviction and KV cache retrieval strategies. KV eviction strategies reduce the size of the KV cache by systematically discarding KV pairs based on predefined policies~\cite{streamingllm,h2o,snapkv,pyramidkv}. However, this usually suffers from significant accuracy loss, as the removal of essential information without comprehensive contextual understanding adversely affects task performance. In contrast, KV retrieval methods maintain accuracy by preserving all KV pairs in the cache and selectively retrieving relevant tokens during the decoding stage, thereby ensuring that no crucial information is omitted~\cite{quest,shadowkv}. Nevertheless, KV retrieval does not mitigate high memory usage, limiting scalable batch sizes and thus the throughput under long contexts. Furthermore, certain approaches~\cite{minference,duoattention,razorattention} recognize the heterogeneous sparsity patterns across different attention heads within the model and configure the heads with distinct sparsification strategies during inference.
While they show initial promise in accuracy and efficiency, they do not fully utilize the interplay between the pre-filling and decoding phases to facilitate further optimization.
% Additionally, selecting top‑k elements during decoding introduces extra computational overhead and approximation errors within the attention mechanism, causing minor latency increases and accuracy reductions.

Different from previous methods, we explore leveraging the experience in pre-filling stage to promote the decoding stage. Our intuition is inspired by an interesting observation: certain attention heads exhibit sequentially consistent attention patterns, spanning across the pre-filling and decoding stages, while some attention heads exhibit rich activation dynamics during the whole process, as shown in Figure~\ref{fig:sequential_consistency}.
Besides, this phenomenon, reflecting the special head working mechanism, frequently existed in different LLMs and their different layers. 
This indicates that if we can capture those modes at the pre-filling stage, we can leverage the sequential consistency to only maintain a subset of crucial tokens for certain heads, which achieves both a reduction of the KV cache and the speedup of attention computation during the decoding stage for LLM inference acceleration. 
% Therefore, identifying these consistent heads is key to fully exploiting this phenomenon.


% Our method employs an observation window during the pre-filling stage to identify each attention head's pattern and its corresponding critical tokens. Specifically, we classify attention heads using the Variable Tracking task dataset from the RULER benchmark as the reference dataset. Within the observation window, we record the frequency of tokens designated as critical and calculate the coefficient of variation to distinguish between consistent and adaptive heads. Specifically, we allocate minimal cache budgets to consistent heads and substantial cache budgets to adaptive heads to cache the top‑k most frequent tokens.
Based on the above analysis, we propose CateKV, a simple, effective, and plug-and-play method designed to enhance LLM inference efficiency by leveraging sequential consistency. CateKV uses an observation window during the pre-filling stage to identify critical tokens and employs a coefficient-of-variation-based score to classify attention heads into consistent and adaptive types, based on a reference dataset. In consistent heads, CateKV retains only a small proportion of critical KV pairs, while in adaptive heads, it retains most, selecting tokens based on their importance derived from attention weights in the pre-filling stage, as shown in Figure~\ref{fig:intro1}. 
By such an effective routing manner guided by sequential consistency, our method maintains the performance merit of LLM inference and simultaneously enjoys the acceleration gain.
Our method is orthogonal to and combinable with many existing acceleration approaches, and we conducted extensive experiments on widely used benchmarks including RULER~\cite{ruler}, Longbench~\cite{longbench}, and NIAH~\cite{niah}, using models such as LLaMA-3-8B-Instruct-1048K~\cite{llama3-8b}, GLM-4-9B-1M~\cite{glm}, LLaMA-3.1-8B~\cite{llama3-1} and Yi-9B-200K~\cite{yi} to demonstrate the effectiveness. 
% The results show that our method achieves performance comparable to full attention, with an $2.14\times$ reduction in memory usage and an $1.81\times$ speedup in decoding on a single A100 GPU.
In a nutshell, our contributions are summarized as follows:
\begin{itemize}
%[itemsep=1pt,parsep=1pt,topsep=1pt,partopsep=1pt]
% v1
% \item An observation of sequential consistency in attention heads, leading to their classification into consistent and adaptive heads.
% \item CateKV, an acceleration method that uses a coefficient-of-variation-based score to classify attention heads, achieving an optimal balance between efficiency and accuracy.
% % An new algorithm 
% \item A comprehensive evaluation of CateKV, demonstrating its ability to reduce memory usage by $2.14\times$ and accelerate decoding by $1.81\times$, while preserving high accuracy.
% v2
% \item We observe a sequential consistency phenomenon in attention heads, and propose a coefficient-of-variation-based score to classify them.
% \item We introduce CateKV, a hybrid KV cache acceleration algorithm based on sequential consistency.
% % An new algorithm 
% \item CateKV reduces memory usage by $2.14\times$ and decoding latency by $1.81\times$, while maintaining accuracy comparable to full attention on long-context benchmarks.
% v3
% \item We observed the phenomenon of sequential consistency in attention heads and further discovered that different heads exhibit varying degrees of consistency, naturally dividing them into consistent and adaptive heads.
\item We identify that general sequential consistency exhibits in certain attention heads and dynamic activation in others, dividing the attention heads into consistent and adaptive heads, which naturally constructs the basic for decoding acceleration with pre-filling experience.
\item We propose CateKV, a hybrid KV cache acceleration algorithm leveraging sequential consistency. 
Employing a coefficient-of-variation-based score, \shortname can precisely classify attention heads into two categories, enabling efficient KV cache eviction while closely approximating the performance of full attention.
% while maintaining the dynamic capability of full attention.
% while讲adaptive head能保留类似full attention的能力（所以第三点能preserving performance），这个能力到底咋写
\item Extensive evaluations on popular benchmarks demonstrate that CateKV reduces memory usage and decoding latency by $2.72\times$ and $2.18\times$ for single inputs, increases throughput by $3.96\times$ in batch scenarios, while maintaining performance comparable to full attention.
% Moreover, our method is plug-and-play, allowing easy integration with other approaches for further acceleration.
Further acceleration can be achieved by integrating our plug-and-play method with other approaches.


\end{itemize}

\section{Related Work}
\subsection{KV Cache Eviction Algorithm}
% 在long-context的场景下，KV cache会随着输入长度的增加而增加，这会导致巨大的时间和空间开销。一些已有的方法通过维护一个固定大小的KV cache来保存一些关键的tokens，并驱逐剩余不重要的tokens来减少内存使用且减少计算量。

To save the significant time and space overhead as the input length increases, various approaches explore evicting tokens to reduce both memory usage and computational cost. StreamingLLM~\cite{streamingllm} introduces the phenomenon of "attention sink" and supports longer sequence by retaining only the KV pairs of attention sinks and recent tokens. H2O~\cite{h2o} employs a low-cost eviction strategy to maintain a fixed-size KV cache containing heavy-hitters, based on the sum of historical attention scores. SnapKV~\cite{snapkv} uses the last tokens in the prompt during the prefilling stage to select critical tokens for generation in the decoding stage. PyramidKV~\cite{pyramidkv}, PyramidInfer~\cite{pyramidinfer}, and LazyLLM~\cite{lazyllm} leverage attention distribution across layers to dynamically adjust cache size, 
% allocating more to lower layers and less to higher layers, thereby 
making cache allocation more efficient. 
Other methods like MagicPig~\cite{magicpig}, Q-Hitter~\cite{q-hitter}, ALISA~\cite{alisa} and , which combine KV cache eviction with quantization, hashing algorithms, or sparse windows, can also improve inference efficiency.
% However, all of these methods lead to varying degrees of performance degradation, as they permanently evict certain tokens that may play a crucial role in future generation processes.
However, these methods induce non-negligible performance degradation, as they potentially evict certain tokens that are crucial for future generation.

\subsection{KV Cache Retrieval Algorithm}
Unlike KV cache eviction algorithms, KV cache retrieval algorithms retain a complete KV cache and dynamically retrieve important KV pairs to reduce inference latency. 
% Quest~\cite{quest} uses PageAttention~\cite{vllm}, dividing tokens into pages, and applying an approximate attention score to retrieve the most relevant pages for the current decoding steps. 
Following PageAttention~\cite{vllm}, Quest~\cite{quest} divides tokens into pages and devises an approximate attention score to retrieve the most relevant pages for the current decoding steps.
InfLLM~\cite{infllm} adopts a strategy similar to Quest, offloading most of the cache to the CPU to support longer prompts.
% ShadowKV~\cite{shadowkv} enhances the storage efficiency of Quest by storing only low-rank key caches and offloading value caches, reducing memory usage and enabling inference with large batch sizes and long contexts. 
ShadowKV~\cite{shadowkv} enhances the storage efficiency of Quest by storing only low-rank key caches and offloading value caches, allowing inference with large batch sizes and context lengths.
Others like SparQ~\cite{sparq}, InfiniGen~\cite{infinigen}, and Loki~\cite{loki}, accelerate the selection of top-\textit{k} critical tokens by reducing the dimension. 
KV cache retrieval methods maintain performance by preserving the entire KV cache but inevitably incur increased inference latency and storage costs.

\subsection{Head-wise Attention Classification}
% Another line of work focuses on classifying attention heads based on their characteristics, with the goal of designing more efficient algorithms to reduce inference latency. 
Another line of work classifies attention heads into distinct sparse patterns.
MInference~\cite{minference} divides the attention into A-shape, Vertical-Slash, and Block-Sparse patterns, achieving acceleration during the \emph{pre-filling} stage. RazorAttention~\cite{razorattention} and DuoAttention~\cite{duoattention} split heads into retrieval heads and streaming heads to determine whether to implement full attention or an attention mechanism similar to StreamingLLM~\cite{streamingllm}. Methods like AdaKV~\cite{adakv} and HeadKV~\cite{headkv} achieve more fine-grained classification by allocating different budgets to each attention head.
These methods focus on the features within individual heads or their variations, but overlook the patterns of attention heads across the prefilling and decoding stages.
\section{Method}

\begin{figure}[t]
\begin{center}
    \includegraphics[width=0.48\textwidth]{Figures/sequential_consistency.pdf}
    \caption{Visualization of attention weight heatmaps for consistent heads and adaptive heads in the LongChat-7B model using a randomly selected text segment. In these figures, the vertical axis represents the attention of the head across different queries, with the first 20 rows corresponding to the last 20 queries during the pre-filling stage, while the subsequent rows depict the attention weights across consecutive decoding steps. These two types of heads exist in layers and popular models (Appendix~\ref{appendix-consistency}).}
    \vspace{-0.5cm}
    \label{fig:sequential_consistency}
\end{center}
\end{figure}

% To effectively explore the consistency of individual attention heads across consecutive decoding steps and their relationship with the pre-filling stage, we randomly selected a text segment from wikitext~\cite{wikitext} as input for the LongChat-7B ~\cite{longchat} model to observe the attention patterns during decoding.

% Numerous related works~\cite{snapkv,minference} have utilized the final few tokens as an observation window to obtain global information during the pre-filling stage. To facilitate subsequent discussions, we denote the size of this observation window as $L_{obs}$ and refer to the attention matrix based on the observation window as the observation matrix. 
% Figure~\ref{fig:sequential_consistency} illustrates the attention heat map based on the observation window ($L_{obs}=10$) during the pre-filling stage, as well as the attention heat maps across all decoding steps for two types of attention heads. 

% We observe that for some attention heads, the attention is concentrated on a few critical tokens, which exhibit clear consistency across various decoding steps. Moreover, these critical tokens can be identified through the attention weights within the observation window during the pre-filling stage. This consistency allows us to guide the decoding process using information derived from the pre-filling stage.

% In contrast, other attention heads exhibit attention weight distributions that vary significantly across decoding steps, maintaining a consistently broader attention scope and placing less emphasis on specific tokens at each step. For these heads, it is essential to retain most of the KV pairs to ensure prediction accuracy. 
% Additionally, we can leverage the sequential consistency in their attention patterns to reduce part of the attention search space, focusing the attention on more relevant and useful information.
\subsection{Sequential Consistency of Attention Heads}
\label{seq_consistency}

In this paragraph, we present an interesting observation about the attention patterns across the pre-filling and decoding stages. To illustrate this, we randomly selected a text segment from Wikitext~\cite{wikitext} as input for the LongChat-7B~\cite{longchat} model, examining how attention weights evolve throughout the generation process.
Figure~\ref{fig:sequential_consistency} illustrates the attention heatmaps during the pre-filling and decoding stages for two types of attention heads. We observe that for certain attention heads, attention is concentrated on a few critical tokens, which show clear consistency across both the pre-filling and decoding stages. This consistency allows us to identify critical tokens during the pre-filling stage, which can then guide the decoding process and help reduce computational costs. In contrast, other attention heads exhibit attention distributions that vary significantly across decoding steps, maintaining a broader attention scope without focusing on specific tokens at each step. For these heads, it is crucial to retain most of the KV pairs to ensure accurate predictions. For clarity, we provide the further claims for these two types of attention heads:

% \begin{itemize}
%   \item \textbf{Consistent heads} are attention heads that focus on a consistent and limited number of critical tokens across decoding steps.
%   \item \textbf{Adaptive heads} are attention heads with variable attention weights across different decoding steps, and attempting to apply consistent patterns to them would negatively affect the model's output.
% \end{itemize}

\begin{itemize}
  \item \textbf{Consistent heads} are attention heads that exhibit stable attention patterns of sequential consistency, focusing on a limited set of tokens across all decoding steps.
  \item \textbf{Adaptive heads} are attention heads characterized by variable attention distributions across decoding steps, which do not exhibit stable patterns and require a larger attention space for flexible token interactions.
\end{itemize}

\subsection{How to Identify Consistent and Adaptive Heads?}
\label{identification_method}
% Building upon our observations, the attention patterns of consistent heads can be identified in the pre-filling stage. To accurately classify consistent and adaptive heads, we propose three distinct, theoretically grounded scores derived from the observation matrix at this stage.

% \begin{figure}[t]
% \begin{center}
%     \includegraphics[width=0.45\textwidth]{Figures/sample_similarity.pdf}
%     \caption{Head Classification Results for 5 Random Samples in Wikitext103 using the Dynamic Identification Method on the  LLaMA-3-8B-Instruct-1048K Model. Each row represents the classification of 8 heads in a layer, with a total of 32 layers. The proportion of consistent heads is 60\%. From the figure, we can clearly observe significant similarity in head classification results across different samples.\jhy{(jhy:Placeholder figure; alternative methods are needed to visualize similarity.}}
%     \label{fig:sample_similarity}
% \end{center}
% \end{figure}

% \paragraph{Sample-Specific Dynamic Identification Method} 
% In the pre-filling stage, a common approach to obtain global information is by observing a small part of the attention weights. In this work, we set an observation window that contains the last portion of the input, with a length of \(L_{\text{obs}}\). Additionally, since initial and recent tokens are typically important but do not influence the identification process, we temporarily exclude them during the identification phase. Their respective lengths are denoted as \(L_{\text{init}}\) and \(L_{\text{rec}}\).

\paragraph{Observation matrix} 
% In the pre-filling stage, a common approach~\cite{snapkv,minference} to obtain global information is by observing a small part of the attention weights. Our observations presented in Section~\ref{seq_consistency} also substantiate this point. In this work, we set an observation window that contains the last portion of the input tokens, with a length of \(L_{\text{obs}}\). Additionally, since initial and recent tokens are typically important, but do not influence the identification process, we temporarily exclude them during the identification phase. Their respective lengths are denoted as \(L_{\text{init}}\) and \(L_{\text{rec}}\). 

As can be seen in Figure~\ref{fig:sequential_consistency}, the attention weights of the last query tokens at the pre-filling stage effectively reflect the attention patterns. To save the computational cost, we set an observation window that contains the last query tokens of the input to identify head types and critical tokens.
% This approach has also been proven effective in similar works~\cite{snapkv,minference}. 
Additionally, since initial and recent tokens are typically important but do not influence the classification process, we temporarily exclude them during the identification phase. Formally, let \(L_{\text{init}}\), \(L_{\text{rec}}\) and \(L_{\text{obs}}\) respectively denote the lengths of the initial tokens, recent tokens, and the observation window. Then, consider a head within a sample at the pre-filling stage, where the input includes the query \( Q\), key \(K\), and value \(V \in \mathbb{R}^{n \times d} \), with \( n \) representing the input length and \( d \) representing the head dimension in the attention mechanism. We define the observation matrix $A$ within the observation window, as follows:
\begin{align}\label{eq:observation-matrix}
    A = \text{softmax}\left(\frac{Q_{[-L_{\text{obs}}:, :]} K_{[L_{\text{init}}:-L_{\text{rec}}, :]}^T}{\sqrt{d}}\right),
\end{align}
where $A\in \mathbb{R}^{L_{\text{obs}} \times \left(n-L_{\text{rec}}-L_{\text{init}}\right)}$.
In the following, we will identify consistent heads based on the characteristics of $A$.

% For clarity, We also define \( L_{c}=n-L_{init}-L_{rec} \) to represent the length of the context window.

% \paragraph{Cosine Similarity(CS) Score}
% A fundamental characteristic of a consistent head is the high degree of attention consistency across sequences, which can be effectively quantified using cosine similarity. To mitigate the influence of atypical attention distributions (such as instances where attention is concentrated on a very limited number of tokens) on the consistency evaluation, we first binarize the attention for each token within the observation window using a percentile-based threshold \( k \). 
% \[
% B_{i,j} = \mathbb{I}\left( A_{i,j} \geq a_{(k,i)} \right)
% \]
% where \( \mathbb{I} \) is the indicator function, and \( a_{(k,i)} = \text{Quantile}_{k}(A_{i,\cdot}) \) represents the \( k \)-th percentile element in the \( i \)-th row of matrix \( A \).

% Subsequently, we calculate the cosine similarity between these binarized attention vectors and the mean attention vector. By averaging these similarity values, we derive the cosine similarity score for this head, denoted as $score_{cos}$.

% \[
% u_{j} = \mathbb{I}\left( \frac{1}{L_{obs}} \sum_{i=1}^{L_{obs}} A_{i,j} \geq a_{k} \right)
% \]

% where \( a_{k} \) is the \( k \)-th element of the mean attention vector.

% \[
% score_{CS} = \frac{1}{m} \sum_{i=1}^{m} \frac{B_{i} \cdot u}{\|B_{i}\| \, \|u\|}
% \]

% \paragraph{Non-Critical Token Ratio(NCTR) Score}
% In a consistent attention head, the interactions between different tokens exhibit not only mutual consistency but also a concentration on a limited subset of critical tokens. To evaluate whether a reduced search space can satisfy the internal attention requirements of the head, we quantify the number of tokens within the context window that are not designated as critical. Similar to cosine similarity–based identification methods, we first binarize the observation matrix and then compute the non-critical token ratio as a score.
% % \[
% % s_{j} = \sum_{i=1}^{L_{obs}} B_{i,j}
% % \]
% % \[
% % score_{NCTR} = \frac{1}{L_{c}} \sum_{j=1}^{L_c} \mathbb{I}\left( s_j = 0 \right)
% % \]
% \[
% score_{NCTR} = \frac{1}{L_{c}} \sum_{j=1}^{L_c} \mathbb{I}\left( \sum_{i=1}^{L_{obs}} B_{i,j} = 0 \right)
% \]

\begin{figure}[t]
\begin{center}
\includegraphics[width=0.48\textwidth]{Figures/static.pdf}
    \caption{(a) High similarity of classification results across samples. We calculate the overlap rate between the classification results of 200 randomly selected WikiText samples and the overall classification result with $r=0.3$.  All samples exhibit an overlap rate exceeding 80\% with the overall result. (b) Frequency of adaptive head identification for Llama-3B model in a reference dataset with $r=0.3$. The frequency distribution is highly concentrated, enabling the determination of a fixed head type based on these results. }
    \vspace{-0.5cm}
    \label{fig:static}
\end{center}
\end{figure}

\paragraph{Coefficient of Variation (CV) Score}
% To simultaneously evaluate the consistency and concentration of attention within the observation matrix, we propose a coefficient of variation score. Since the coefficient of variation is highly sensitive to the magnitude of values, we also binarize the observation matrix first. We continue to employ a percentile-based threshold \( k\) to reduce the influence of the mean and introduce a scaling factor \( \alpha \) to prevent the omission of values near the threshold.
Empirically, consistent heads exhibit two primary features in their observation matrices: a high degree of similarity among different rows and a small subset of columns that is sufficient to recall most of the attention weights. Inspired by this aspect, we propose a coefficient-of-variation-based score to measure the concentration and similarity of attention within observation matrices $A$. Since the coefficient of variation~\cite{cv} is highly sensitive to the magnitude of values, we first binarize the observation matrix with a percentile-based threshold \( k\) and a scaling factor \( \alpha \).
\begin{align}\label{eq:binarization}
    B = \mathbb{I}\left( A \geq \Phi(k,\alpha) \right) \in \mathbb{R}^{L_{\text{obs}} \times \left(n-L_{\text{rec}}-L_{\text{init}}\right)}
\end{align}
% with \( a_{(k,\alpha)}=\text{Quantile}_{k}(A)\ \times \alpha \)
where \( \mathbb{I} \) is the indicator function, and \( \Phi(k,\alpha) = \text{Quantile}_{k}(A) \times \alpha \) represents the \( k \)-th percentile element in the matrix \( A \), scaled by a factor of \( \alpha \). Then after binarization, we derive a frequency vector 
\( C \) to quantify the number of times each token that is identified as critical, reflecting the similarity and concentration of attention weights:
\begin{align}\label{eq:frequency-vector}
C = \sum_{i=0}^{L_{\text{obs}}} B_{i,:}\in \mathbb{R}^{\left(n-L_{\text{rec}}-L_{\text{init}}\right)} 
\end{align}
Now, we can obtain the score based on the coefficient of variation for an attention head of a sample as follows:
\begin{align}\label{eq:cv-score}
\text{score} = \frac{\sqrt{\frac{1}{\left(n-L_{\text{rec}}-L_{\text{init}}\right)} \sum_{i} (C_i - \mu(C))^2}}{\mu(C)} 
\end{align}
where \(\mu(C) = \frac{1}{\left(n-L_{\text{rec}}-L_{\text{init}}\right)} \sum_{i} C_i\) is the mean. With the above equations, we can compute a score for each head under a sample. Then, for a specific LLM, we can obtain a score matrix \( S \in \mathbb{R}^{l \times h}\) for all heads, where \(l\) and \(h\) represent the number of layers in the model and the number of heads in a layer, respectively. For a statistic head identification rule, let $r$ denote the proportion of adaptive heads, and \( \Gamma(r) \) represents the percentile threshold based on \( r \). We can use the threshold to control the token eviction ratio. Then, we distinguish the head types of a specific sample as follows
\begin{align}\label{eq:head-classification}
\text{Head}_{i,j} = 
\begin{cases} 
\text{consistent head}, & \text{if } S_{i,j} > \Gamma(r) \\
\text{adaptive head}, & \text{if } S_{i,j} \leq \Gamma(r) 
\end{cases}.
\end{align}

% Comparing these three scores, we consider that the CV score is the most effective because a higher coefficient of variation indicates that frequencies are concentrated on the same tokens, reflecting both consistency and attention aggregation. In Section~\ref{different_methods_ablation}, we validate this by demonstrating the performance of methods based on these three scores. In subsequent experiments, we use the CV score to identify consistent heads.

% Next, we define a threshold to determine which positions are relatively important for the current attention. The threshold is given by 
% \[ 
% \text{threshold} = v_{\beta} \times \alpha \]
% where \( v_{\beta} \) is the value at the \( \beta \)-th percentile of \( A_{\text{obs}} \), and \( \alpha \) is a scaling factor. Here, \( \beta \) is a hyperparameter referred to as the high-attention percentile.

% We can obtain a binary matrix \( A_{\text{obs}}' \) and a frequency vector \( c \) as follows:
% \[
% A_{\text{obs}}' = A_{\text{obs}} > \text{threshold} 
% \]
% \[
% c = \sum_{k=0}^{L_{obs}} A_{\text{obs}}'[k,:]
% \] 
% Now, we can measure the concentration of attention within the head by calculating the coefficient of variation (CV) of \( c \):
% \[
% \text{CV}(c) = \frac{\sqrt{\frac{1}{n} \sum_{k=1}^{n} (c_k - \mu(c))^2}}{\mu(c)}
% \]
% where \(\mu(c) = \frac{1}{n} \sum_{k=1}^{n} c_k\)

\begin{algorithm}[tb]
   \caption{CateKV in an individual Head}
   \label{alg:method}
\begin{algorithmic}
   \STATE {\bfseries Input:} $Q, K, V \in \mathbb{R}^{n \times d}$, $q \in \mathbb{R}^{1 \times d}$, head type $H$, observation window size $L_{\text{obs}}$, selected chunk budget $b$, chunk size $cs$, retention ratio $\eta$
   \STATE {\bfseries Pre-filling Stage:} 
   
   \LineComment{Calculating token criticality}
   
    \STATE $A = \text{softmax}(Q_{[-L_{\text{obs}}:, :]} K_{[:-L_{\text{obs}}, :]}^T / {\sqrt{d}} )$
   \STATE $C = \sum_{i=0}^{L_{\text{obs}}} A_{i,:}$

   \LineComment{Divide $C$ into chunks and take maximum}
   \STATE $C_{\text{chunk}} = \text{BlockMax}(C,cs)$

   \LineComment{Cache keys and values based on the indices of top-k elements}
   
   \IF{$H=\text{consistent head}$}  
   \STATE $i_k = \text{argtopk}(C_{\text{chunk}},b)$
   \STATE $cache_\text{k}, cache_\text{v} = \text{Cache}(K,V,i_k,L_{\text{obs}})$
   \ELSE
   \STATE $i_k = \text{argtopk}(C_{\text{chunk}},n\eta)$
   \STATE $cache_\text{k}, cache_\text{v} = \text{Cache}(K,V,i_k,L_{\text{obs}})$
   \ENDIF

   \STATE {\bfseries Decoding Stage:}

   \LineComment{Retrieval keys and values from cache}
   
   \IF{$H=\text{consistent head}$} 
   \STATE $k,v = cache_\text{k}, cache_\text{v}$
   \ELSE
   \STATE $k,v = \text{Retrival}(q, cache_\text{k}, cache_\text{v})$
   \LineComment{all or query-aware}
   
   \ENDIF
   \STATE $output = \text{Attention}(q, k, v)$ 
\end{algorithmic}
% \vspace{-0.2cm}
\end{algorithm}

\paragraph{Reference-Based Static Identification}
Although head identification can be performed dynamically for each sample, it is actually expensive for memory management along with the change of samples at the pre-filling stage. Therefore, it is better to determine a fixed head type for the model, which can be directly used for inference. Surprisingly, we observe that the same head does exhibit similar attention patterns across different samples, and the identification results based on CV scores are highly consistent. This inspires us to use a reference dataset to calculate the frequency of each head being identified as the consistent head, and then derive the final model-wise identification result based on the adaptive head ratio $r$, as shown in Figure~\ref{fig:static}, which comes to the final form of our method (termed as \emph{CateKV}). After determining the type of each head, we retain only the most important KV pairs for consistent heads, while for adaptive heads, we preserve a majority based on the predefined retention ratio \( \eta \). Specifically,
we select the top-$k$ KV pairs in chunks, enabling seamless integration with retrieval-based methods. For the GQA model, the observation matrix \( A \) is computed as the mean of \( A \) of the heads in a group. 
% In adaptive heads, we employ the Quest as the KV retrieval method during the decoding stage here. 
The CateKV acceleration for LLMs is shown in Algorithm~\ref{alg:method}.
% Under the same identification method, we find that the classification results across samples are very similar. Thus, we can use a reference dataset to calculate the frequency of each head being identified as the consistent head, and then derive the final model-wise classification result based on the consistent head ratio \( r \).
% After calculating the scores for each head in each layer, we perform binary classification of head \( j \) in layer \( i \) using a sparse percentile threshold \( \tau_d \):

% \paragraph{Similarity-Based Static Identification Method}
% Although the dynamic identification method allows for determining the head type for each sample individually, it introduces additional overhead during the pre-filling stage of inference. Moreover, because the dynamic identification method requires hyperparameters to determine the threshold, it is not very convenient for deployment. We observed that different samples consistently exhibit similar results in head classification. Figure~\ref{fig:sample_similarity} shows the classification results for 5 random samples in wikitext-103 under the dynamic method. Based on this observation, we propose a similarity-based static identification method to get model-wise head classification results.

% In the static method, we use Wikitext-103 as our reference dataset. We concatenate all the samples in reference dataset and then perform truncation to construct our dataset, with the dataset size set to \( M \). We then use the dynamic identification method on all samples to obtain the frequency of consistent heads. Finally, we apply a sparse percentile threshold \( \tau_s \) to perform the final classification.
% \[
% H = \frac{1}{M} \sum_{k=1}^{M} H_k
% \]

% \[
% \text{head\_s}_{i,j} = 
% \begin{cases} 
% \text{consistent head}, & \text{if } \text{H}_{i,j} > \tau_s \\
% \text{adaptive head}, & \text{if } 
% \text{H}_{i,j} \leq \tau_s 
% \end{cases}
% \]
% Here, \( H \in \mathbb{R}^{l \times h} \) represents the dynamic identification results, where \( l \) is the number of layers and \( h \) is the number of attention heads per layer. \( H \) is a binary matrix, where 1 indicates that the corresponding position is identified as a consistent head.


\begin{figure}[t]
\begin{center}
    \includegraphics[width=0.42\textwidth]{Figures/4method.pdf}
    \caption{Attention weight recall curves for four methods. The curves show the attention recall for each attention head, calculated on four patterns: consistent, streaming, quest, and top-k, on a 128k-length example. The attention recall values are obtained by the mean of all decoding steps and sorted in descending order for all heads. The sparse budget is set to 2048, and the chunk size is 8. The results indicate that the consistent pattern outperforms the streaming pattern in terms of overall attention recall, and approaches the performance of the quest pattern that requires additional computation.}
    \vspace{-0.3cm}
    \label{fig:4method}
\end{center}
\end{figure}

% \subsection{Relations with KV Eviction and Retrieval Methods}
% As shown in Figure~\ref{fig:intro1}, the KV Eviction method assumes that the criticality of the tokens can be directly derived from prior knowledge in the pre-filling stage and remains useful throughout the decoding stage, while the KV Retrieval method suggests that the criticality can change with different query tokens. Both methods focus on a particular attention pattern within the model, but neither of them is always accurate. By further exploring the model's internal characteristics, we find that consistent heads and adaptive heads exhibit these two patterns, respectively, and we can distinguish them using a reference-based approach. Therefore, we can easily combine these two methods to leverage their advantages in efficiency and accuracy. A detailed discussion can be found in Section~\ref{combination}.

% \subsection{Consistent Head vs. Streaming Head}
% Related works~\cite{razorattention, duoattention} classify attention heads into Retrieval Heads and Streaming Heads, which is similar to our method. In this section, we discuss the relationship between these two classifications and the advantages of our method.
% \begin{table*}[t]
\centering
\caption{Performance (\%) of different models and various methods on RULER evaluated at length of 128K. The 'Cache' in the table refers to the retained KV cache size. \shortname outperforms other methods and comparable with full attention. }
\setlength{\tabcolsep}{3pt} % 调整列间距
\small

\begin{tabular}{l|c|cccccccccccc|c}
\toprule
Methods & Cache & N-S1 & N-S2 & N-S3 & N-MK1 & N-MK2 & N-MK3 & FWE & N-MQ & N-MV & QA-1 & QA-2 & VT & Avg. \\ 
\midrule
\textit{Llama-3-8B-1M} & 100\% & 100.00 & 100.00 & 100.00 & 98.96 & 98.96 & 41.67 & 71.88 & 98.69 & 96.35 & 73.96 & 50.00 & 78.75 & 84.10 \\
StreamingLLM  & 41\% & 51.04 & 51.04 & 51.04 & 40.63 & 35.42 & 22.92 & 75.69 & 45.31 & 39.58 & 78.13 & 45.83 & 31.46 & 47.34 \\
SnapKV & 41\% & 100.00 & 100.00 & 100.00 & 98.96 & 98.96 & 30.21 & 71.53 & 98.44 & 97.13 & 73.95 & 51.04 & 79.17 & 83.28 \\
PyramidKV  & 41\% &100.00 &	100.00 &	100.00 & 98.96 & 98.96 & 37.50 & 71.53 &	98.44 &	96.61 &	71.87 &	50.00 &	79.38 & 83.60 \\
Duoattention & 41\% & 100.00 &	100.00 &	100.00 &	98.96 &	97.92 &	39.58 	&76.74& 	94.27 	&90.36 &	69.79 &	51.04 &	86.46 &	83.76 \\
\rowcolor{cyan!10}
CateKV  & 41\% & 100.00 & 100.00 & 100.00 & 98.96 & 97.92 & 41.67 & 71.88 & 98.44 & 96.88 & 73.96 & 50.00 & 85.63 & 84.61 \\
\midrule
\textit{Phi-3-Mini-128K} & 100\% & 96.88 & 90.63 & 95.83 & 83.33 & 65.63 & 37.50 & 87.15 & 72.14 & 66.67 & 63.54 & 39.58 & 65.83 & 72.06 \\
StreamingLLM & 41\% & 47.91 & 45.83 & 44.79 & 38.54 & 30.21 & 25.00 & 84.38 & 36.49 & 34.90 & 64.58 & 38.54 & 4.38 & 41.29 \\
SnapKV  & 41\% & 96.88 & 90.63 & 80.21	& 82.29 & 56.25 & 11.46 & 82.99 &	61.72 & 53.91 &	62.50 & 38.54 & 63.54 & 65.08 \\
PyramidKV  & 41\% & 96.88 &	90.63 &	84.38 &	83.33 &	57.29 &	13.54 &	78.47 &	66.15 &	59.64 &	62.50 &	39.58 & 62.29 & 66.22 \\
\rowcolor{cyan!10}
CateKV  & 41\% & 96.88 & 90.63 & 95.83 & 83.33 & 65.63 & 38.54 & 80.21 & 70.31 & 65.63 & 63.54 & 39.58 & 70.21 & 71.69 \\
\midrule
\textit{Llama-3.1-8B} & 100\% & 100.00 & 100.00 & 98.96 & 98.96 & 90.63 & 63.54 & 71.53 & 98.96 & 95.31 & 81.25 & 46.88 & 68.54 & 84.55 \\
StreamingLLM  & 41\% & 51.04 & 51.04 & 51.04 & 39.58 & 34.38 & 40.63 & 71.18 & 44.27 & 39.84 & 85.42 & 40.63 & 28.33 & 48.11 \\
SnapKV   & 41\% & 100.00 & 100.00 & 98.96 & 98.96 & 89.58 & 46.88 & 69.10 & 98.96 & 94.01 & 81.25 & 46.88 & 68.96 & 82.80 \\
PyramidKV  & 41\% & 100.00 & 100.00 & 98.96 & 98.96 & 90.63 &	56.25 &	65.28 &	98.96 &	95.31 &	80.21 &	46.88 & 65.42 & 83.07 \\
Duoattention & 41\% &	100.00 	&100.00 &	98.96 &	97.92 &	88.54 &	59.38 &	74.31 &	97.92 &	91.41 &	81.25 &	46.88 &	78.54 &	84.59 \\ 
\rowcolor{cyan!10}
CateKV  & 41\% & 100.00 & 100.00 & 98.96 & 98.96 & 88.54 & 61.46 & 71.88 & 98.96 & 94.27 & 81.25 & 46.88 & 74.79 & 84.66 \\
\midrule
\textit{Yi-9B-200K} & 100\% & 100.00 &	100.00 &	98.96 &	85.42 &	63.54 &	18.75 &	89.24 	&66.41 &	32.55 &	45.83 &	38.54 &	35.00 &	64.52 \\
StreamingLLM  & 41\% & 47.92 &	52.08 &	50.00 &	39.58 &	37.50 &	7.29 &	90.28 &	33.33 &	14.84 &	44.79 &	36.46 &	18.13 &	39.35 \\
SnapKV  & 41\% & 100.00 &	100.00 	&95.83 &	85.41 &	43.75 &	3.13 &	83.33 &	66.41 &	33.33 &	46.88 &	40.63 &	38.96 &	61.47 \\
PyramidKV  & 41\% & 100.00 &	100.00 &	91.67 &	86.46 &	50.00 &	2.08 &	73.61 &	68.75 &	34.64 &	43.75 &	37.50 &	41.46 &	60.83 \\
\rowcolor{cyan!10}
CateKV  & 41\% & 100.00 &	100.00 	&100.00 &	84.38 	&70.83 &	18.75 &	92.01 &	62.24 &	34.64 &	43.75 &	37.50 &	45.00 	&65.76 \\
\bottomrule
\end{tabular}
\label{tab:ruler-128k-eviction}
\vspace{-0.3cm}
\end{table*}

% % \paragraph{Differences and Connections between Streaming Heads and Consistent Heads}
% \paragraph{Differences}
% Both Consistent Heads and Streaming Heads exploit distinctive attention patterns within the model. However, consistent Heads emphasize the overall similarity of attention across different sequences, while Streaming Heads prioritize whether initial and recent tokens cover most of the attention. In the process of identifying consistent Heads, CateKV considers the initial and recent tokens to be consistently important, classifying based on the remaining tokens. This is in contrast to the classification approach used for Streaming Heads. The attention map in Figure~\ref{fig:streaming} shows that there is no direct correlation between these two types. 

% \paragraph{Advantages}
% Consistent Heads incorporate additional priors from the pre-filling stage, enhancing their overall precision. In the right panel of Figure~\ref{fig:streaming}, we present attention weight recall curves for all attention heads, comparing the consistent and streaming patterns within the same example. The consistent pattern outperforms the streaming pattern and approaches the attention weight recall levels of the KV Retrieval method, such as Quest~\cite{quest}. As the proportion of streaming heads increases, attention recall declines noticeably. Therefore, the high precision of the consistent pattern can be combined with streaming-based methods to achieve a higher compression ratio.

\subsection{Theory Analysis}
In this part, we present an analysis of the theoretical bound on the eviction performance of CateKV.

\textbf{Lemma 1:} Let $G$ represent the hypothesis class derived from the CV-based function, $F$ denote the real-valued function class induced by the binary cross-entropy loss applied to $G$, and let $N$ denote the sample size of the reference-based dataset. Then, with probability at least $1 - \delta$, the following Rademacher complexity bound holds:

\begin{align}
    \forall f \in F, \, &P_{\text{head}}\left( \mathbb{E}[f] - \frac{1}{N} \sum_{n=1}^{N} f_n \leq 2 \mathcal{R}_N(F) \right) \nonumber \\
    &\quad + \sqrt{\frac{2 \log{\frac{2}{\delta}}}{2N}} \geq 1 - \delta.
\end{align}

where $\mathcal{R}_N(F)$ is the conditional Rademacher average.

Let \( P_1 \) and \( P_2 \) denote the probabilities of correctly classified consistent heads and adaptive heads, respectively, while \( \bar{P}_1 \) and \( \bar{P}_2 \) represent the probabilities of misclassified consistent heads and adaptive heads, respectively. It is assumed that \( P_1 + P_2 = P_{\text{head}} \) and \( \bar{P}_1 + \bar{P}_2 = 1 - P_{\text{head}} \). The probability in the above lemma can then be decomposed through a fine-grained analysis as follows.


% \textbf{Theorem 1:} Let $\eta_1$, $\eta_2$ denote the retention ratios of consistent heads and adaptive heads, respectively, and $\eta_1^*$, $\eta_2^*$ denote their optimal retention ratios correspondingly. Define the retention accuracy for different cases $r_{i,j} = \eta_i^* \mathbb{1}[\eta_j > \eta_i^*] + \eta_j (1 - \mathbb{1}[\eta_j > \eta_i^*])$ by comparing the retention budgets with the optimal budgets, and the hypothesis that the query attention score is the best description in order holds with probability $\lambda$. Then, the token retention accuracy has the following inequalities,

\textbf{Theorem 1:} Let \( \eta_1 \) and \( \eta_2 \) denote the retention ratios for consistent heads and adaptive heads, respectively, while \( \eta_1^* \) and \( \eta_2^* \) represent their corresponding optimal retention ratios. Define the retention accuracy for different cases $r_{i,j} = \eta_i^* \mathbb{I}[\eta_j > \eta_i^*] + \eta_j (1 - \mathbb{I}[\eta_j > \eta_i^*])$ by comparing the retention budgets with the optimal budgets. Additionally, assume that the hypothesis asserting the query attention score provides the best description in order with probability \( \lambda \). Then, the token retention accuracy satisfies:

\begin{align}
    P_{\text{token}} &= \lambda \left( r_{1,1} P_1 + r_{2,2} P_2 + r_{2,1} \bar{P}_1 + r_{1,2} \bar{P}_2 \right) \nonumber \\
    &\geq \lambda \left( \min(r_{2,1}, r_{1,2}) + \left[ \min(r_{1,1}, r_{2,2}) \right. \right. \nonumber \\
    &\quad \left. \left. - \min(r_{2,1}, r_{1,2}) \right] P_{\text{head}} \right).
\end{align}

% \[
% P_{\text{token}} = \lambda \left( r_{1,1} P_1 + r_{2,2} P_2 + r_{2,1} \bar{P}_1 + r_{1,2} \bar{P}_2 \right) \geq \lambda \left( \min(r_{2,1}, r_{1,2}) + \left[\min(r_{1,1}, r_{2,2}) - \min(r_{2,1}, r_{1,2})\right] P_{\text{head}} \right).
% \]

\textbf{Remark 1:} From the above theorem, three critical factors influence the worst-case token retention accuracy (i.e., the lower bound):
\begin{itemize}
    \item \( \lambda \): The effectiveness of identifying an efficient measure to characterize token correlation by the score order, with as high a probability of correctness as possible.
    \item Budget control: The ability to appropriately set the retention budget in order to maximize gains by reducing the majority of tokens when heads are correctly classified, while simultaneously mitigating the negative impact when heads are misclassified.
    \item \( P_{\text{head}} \): The accuracy of the CV-based method in classifying the head type during token reduction, which directly impacts the overall performance.
\end{itemize}


% \textbf{Remark 1:} From the above theorem, three critical factors matter with respect to the worst-token retention accuracy (i.e., the lower bound):
% \begin{itemize}
%     \item $\lambda$, whether we can find an effective and efficient measure to characterize the token correlation by score order that is correct with as high a probability as possible.
%     \item Budget control, whether we can set the proper budget that maximizes the gain by reducing most tokens when heads are correctly classified, while simultaneously weakening the negative effect when heads are misclassified.
%     \item $P_{\text{head}}$, whether the CV-based method can classify the head type as accurately as possible during token reduction.
% \end{itemize}


\begin{table*}[t]
\centering
\caption{Performance (\%) of different models and various methods on RULER evaluated at length of 128K. The 'Cache' in the table refers to the retained KV cache size. \shortname outperforms other methods and comparable with full attention. }
\setlength{\tabcolsep}{3pt} % 调整列间距
\small

\begin{tabular}{l|c|cccccccccccc|c}
\toprule
Methods & Cache & N-S1 & N-S2 & N-S3 & N-MK1 & N-MK2 & N-MK3 & FWE & N-MQ & N-MV & QA-1 & QA-2 & VT & Avg. \\ 
\midrule
\textit{Llama-3-8B-1M} & 100\% & 100.00 & 100.00 & 100.00 & 98.96 & 98.96 & 41.67 & 71.88 & 98.69 & 96.35 & 73.96 & 50.00 & 78.75 & 84.10 \\
StreamingLLM  & 41\% & 51.04 & 51.04 & 51.04 & 40.63 & 35.42 & 22.92 & 75.69 & 45.31 & 39.58 & 78.13 & 45.83 & 31.46 & 47.34 \\
SnapKV & 41\% & 100.00 & 100.00 & 100.00 & 98.96 & 98.96 & 30.21 & 71.53 & 98.44 & 97.13 & 73.95 & 51.04 & 79.17 & 83.28 \\
PyramidKV  & 41\% &100.00 &	100.00 &	100.00 & 98.96 & 98.96 & 37.50 & 71.53 &	98.44 &	96.61 &	71.87 &	50.00 &	79.38 & 83.60 \\
Duoattention & 41\% & 100.00 &	100.00 &	100.00 &	98.96 &	97.92 &	39.58 	&76.74& 	94.27 	&90.36 &	69.79 &	51.04 &	86.46 &	83.76 \\
\rowcolor{cyan!10}
CateKV  & 41\% & 100.00 & 100.00 & 100.00 & 98.96 & 97.92 & 41.67 & 71.88 & 98.44 & 96.88 & 73.96 & 50.00 & 85.63 & 84.61 \\
\midrule
\textit{Phi-3-Mini-128K} & 100\% & 96.88 & 90.63 & 95.83 & 83.33 & 65.63 & 37.50 & 87.15 & 72.14 & 66.67 & 63.54 & 39.58 & 65.83 & 72.06 \\
StreamingLLM & 41\% & 47.91 & 45.83 & 44.79 & 38.54 & 30.21 & 25.00 & 84.38 & 36.49 & 34.90 & 64.58 & 38.54 & 4.38 & 41.29 \\
SnapKV  & 41\% & 96.88 & 90.63 & 80.21	& 82.29 & 56.25 & 11.46 & 82.99 &	61.72 & 53.91 &	62.50 & 38.54 & 63.54 & 65.08 \\
PyramidKV  & 41\% & 96.88 &	90.63 &	84.38 &	83.33 &	57.29 &	13.54 &	78.47 &	66.15 &	59.64 &	62.50 &	39.58 & 62.29 & 66.22 \\
\rowcolor{cyan!10}
CateKV  & 41\% & 96.88 & 90.63 & 95.83 & 83.33 & 65.63 & 38.54 & 80.21 & 70.31 & 65.63 & 63.54 & 39.58 & 70.21 & 71.69 \\
\midrule
\textit{Llama-3.1-8B} & 100\% & 100.00 & 100.00 & 98.96 & 98.96 & 90.63 & 63.54 & 71.53 & 98.96 & 95.31 & 81.25 & 46.88 & 68.54 & 84.55 \\
StreamingLLM  & 41\% & 51.04 & 51.04 & 51.04 & 39.58 & 34.38 & 40.63 & 71.18 & 44.27 & 39.84 & 85.42 & 40.63 & 28.33 & 48.11 \\
SnapKV   & 41\% & 100.00 & 100.00 & 98.96 & 98.96 & 89.58 & 46.88 & 69.10 & 98.96 & 94.01 & 81.25 & 46.88 & 68.96 & 82.80 \\
PyramidKV  & 41\% & 100.00 & 100.00 & 98.96 & 98.96 & 90.63 &	56.25 &	65.28 &	98.96 &	95.31 &	80.21 &	46.88 & 65.42 & 83.07 \\
Duoattention & 41\% &	100.00 	&100.00 &	98.96 &	97.92 &	88.54 &	59.38 &	74.31 &	97.92 &	91.41 &	81.25 &	46.88 &	78.54 &	84.59 \\ 
\rowcolor{cyan!10}
CateKV  & 41\% & 100.00 & 100.00 & 98.96 & 98.96 & 88.54 & 61.46 & 71.88 & 98.96 & 94.27 & 81.25 & 46.88 & 74.79 & 84.66 \\
\midrule
\textit{Yi-9B-200K} & 100\% & 100.00 &	100.00 &	98.96 &	85.42 &	63.54 &	18.75 &	89.24 	&66.41 &	32.55 &	45.83 &	38.54 &	35.00 &	64.52 \\
StreamingLLM  & 41\% & 47.92 &	52.08 &	50.00 &	39.58 &	37.50 &	7.29 &	90.28 &	33.33 &	14.84 &	44.79 &	36.46 &	18.13 &	39.35 \\
SnapKV  & 41\% & 100.00 &	100.00 	&95.83 &	85.41 &	43.75 &	3.13 &	83.33 &	66.41 &	33.33 &	46.88 &	40.63 &	38.96 &	61.47 \\
PyramidKV  & 41\% & 100.00 &	100.00 &	91.67 &	86.46 &	50.00 &	2.08 &	73.61 &	68.75 &	34.64 &	43.75 &	37.50 &	41.46 &	60.83 \\
\rowcolor{cyan!10}
CateKV  & 41\% & 100.00 &	100.00 	&100.00 &	84.38 	&70.83 &	18.75 &	92.01 &	62.24 &	34.64 &	43.75 &	37.50 &	45.00 	&65.76 \\
\bottomrule
\end{tabular}
\label{tab:ruler-128k-eviction}
\vspace{-0.3cm}
\end{table*}


\subsection{Further Discussion of CateKV}
% As shown in Figure~\ref{fig:intro1}, the KV eviction method assumes that the criticality of the tokens can be directly derived from prior knowledge at the pre-filling stage and remains useful throughout the decoding stage, while the KV Retrieval method suggests that the criticality can change with different query tokens. Our observations suggest that neither of these approaches is continuously optimal on its own; rather, a combined usage of both methods is more effective. In this section, we offer an alternative perspective to justify the rationale behind our approach.
To show the distinction and effectiveness of attention patterns discovered from the sequential consistency, Figure~\ref{fig:4method} compares attention weight recall for our consistent pattern, streaming pattern~\cite{streamingllm}, Quest pattern~\cite{quest}, and the upper-bound top-k pattern under the same sparse budget. As can be seen, all methods exhibit a gap compared to the upper bound, indicating some information loss with current sparse attention approaches. Therefore, retaining most KV pairs for certain heads is important for maintaining accuracy. On the other hand, the attention recall of the consistent pattern closely approximates that of the Quest pattern, applying the consistent pattern to heads with sequential consistency can help KV retrieval methods reduce memory usage and the cost of selecting critical tokens. Related works~\cite{razorattention, duoattention} classify attention heads into Retrieval Heads and Streaming Heads, which is similar to our method. However, from an attention recall perspective, the streaming pattern is only effective for a small fraction of heads, otherwise deviating significantly from full attention. This highlights that the consistent pattern achieves a higher compression rate than the streaming pattern. We will discuss the comparison and integration with these methods in the experimental section.
% \subsection{Consistent Head vs. Streaming Head}
% Related works~\cite{razorattention, duoattention} classify attention heads into Retrieval Heads and Streaming Heads, which is similar to our method. In this section, we discuss the relationship between these two classifications and the advantages of our method.
% \begin{table*}[t]
\centering
\caption{Performance (\%) of different models and various methods on RULER evaluated at length of 128K. The 'Cache' in the table refers to the retained KV cache size. \shortname outperforms other methods and comparable with full attention. }
\setlength{\tabcolsep}{3pt} % 调整列间距
\small

\begin{tabular}{l|c|cccccccccccc|c}
\toprule
Methods & Cache & N-S1 & N-S2 & N-S3 & N-MK1 & N-MK2 & N-MK3 & FWE & N-MQ & N-MV & QA-1 & QA-2 & VT & Avg. \\ 
\midrule
\textit{Llama-3-8B-1M} & 100\% & 100.00 & 100.00 & 100.00 & 98.96 & 98.96 & 41.67 & 71.88 & 98.69 & 96.35 & 73.96 & 50.00 & 78.75 & 84.10 \\
StreamingLLM  & 41\% & 51.04 & 51.04 & 51.04 & 40.63 & 35.42 & 22.92 & 75.69 & 45.31 & 39.58 & 78.13 & 45.83 & 31.46 & 47.34 \\
SnapKV & 41\% & 100.00 & 100.00 & 100.00 & 98.96 & 98.96 & 30.21 & 71.53 & 98.44 & 97.13 & 73.95 & 51.04 & 79.17 & 83.28 \\
PyramidKV  & 41\% &100.00 &	100.00 &	100.00 & 98.96 & 98.96 & 37.50 & 71.53 &	98.44 &	96.61 &	71.87 &	50.00 &	79.38 & 83.60 \\
Duoattention & 41\% & 100.00 &	100.00 &	100.00 &	98.96 &	97.92 &	39.58 	&76.74& 	94.27 	&90.36 &	69.79 &	51.04 &	86.46 &	83.76 \\
\rowcolor{cyan!10}
CateKV  & 41\% & 100.00 & 100.00 & 100.00 & 98.96 & 97.92 & 41.67 & 71.88 & 98.44 & 96.88 & 73.96 & 50.00 & 85.63 & 84.61 \\
\midrule
\textit{Phi-3-Mini-128K} & 100\% & 96.88 & 90.63 & 95.83 & 83.33 & 65.63 & 37.50 & 87.15 & 72.14 & 66.67 & 63.54 & 39.58 & 65.83 & 72.06 \\
StreamingLLM & 41\% & 47.91 & 45.83 & 44.79 & 38.54 & 30.21 & 25.00 & 84.38 & 36.49 & 34.90 & 64.58 & 38.54 & 4.38 & 41.29 \\
SnapKV  & 41\% & 96.88 & 90.63 & 80.21	& 82.29 & 56.25 & 11.46 & 82.99 &	61.72 & 53.91 &	62.50 & 38.54 & 63.54 & 65.08 \\
PyramidKV  & 41\% & 96.88 &	90.63 &	84.38 &	83.33 &	57.29 &	13.54 &	78.47 &	66.15 &	59.64 &	62.50 &	39.58 & 62.29 & 66.22 \\
\rowcolor{cyan!10}
CateKV  & 41\% & 96.88 & 90.63 & 95.83 & 83.33 & 65.63 & 38.54 & 80.21 & 70.31 & 65.63 & 63.54 & 39.58 & 70.21 & 71.69 \\
\midrule
\textit{Llama-3.1-8B} & 100\% & 100.00 & 100.00 & 98.96 & 98.96 & 90.63 & 63.54 & 71.53 & 98.96 & 95.31 & 81.25 & 46.88 & 68.54 & 84.55 \\
StreamingLLM  & 41\% & 51.04 & 51.04 & 51.04 & 39.58 & 34.38 & 40.63 & 71.18 & 44.27 & 39.84 & 85.42 & 40.63 & 28.33 & 48.11 \\
SnapKV   & 41\% & 100.00 & 100.00 & 98.96 & 98.96 & 89.58 & 46.88 & 69.10 & 98.96 & 94.01 & 81.25 & 46.88 & 68.96 & 82.80 \\
PyramidKV  & 41\% & 100.00 & 100.00 & 98.96 & 98.96 & 90.63 &	56.25 &	65.28 &	98.96 &	95.31 &	80.21 &	46.88 & 65.42 & 83.07 \\
Duoattention & 41\% &	100.00 	&100.00 &	98.96 &	97.92 &	88.54 &	59.38 &	74.31 &	97.92 &	91.41 &	81.25 &	46.88 &	78.54 &	84.59 \\ 
\rowcolor{cyan!10}
CateKV  & 41\% & 100.00 & 100.00 & 98.96 & 98.96 & 88.54 & 61.46 & 71.88 & 98.96 & 94.27 & 81.25 & 46.88 & 74.79 & 84.66 \\
\midrule
\textit{Yi-9B-200K} & 100\% & 100.00 &	100.00 &	98.96 &	85.42 &	63.54 &	18.75 &	89.24 	&66.41 &	32.55 &	45.83 &	38.54 &	35.00 &	64.52 \\
StreamingLLM  & 41\% & 47.92 &	52.08 &	50.00 &	39.58 &	37.50 &	7.29 &	90.28 &	33.33 &	14.84 &	44.79 &	36.46 &	18.13 &	39.35 \\
SnapKV  & 41\% & 100.00 &	100.00 	&95.83 &	85.41 &	43.75 &	3.13 &	83.33 &	66.41 &	33.33 &	46.88 &	40.63 &	38.96 &	61.47 \\
PyramidKV  & 41\% & 100.00 &	100.00 &	91.67 &	86.46 &	50.00 &	2.08 &	73.61 &	68.75 &	34.64 &	43.75 &	37.50 &	41.46 &	60.83 \\
\rowcolor{cyan!10}
CateKV  & 41\% & 100.00 &	100.00 	&100.00 &	84.38 	&70.83 &	18.75 &	92.01 &	62.24 &	34.64 &	43.75 &	37.50 &	45.00 	&65.76 \\
\bottomrule
\end{tabular}
\label{tab:ruler-128k-eviction}
\vspace{-0.3cm}
\end{table*}

% % \paragraph{Differences and Connections between Streaming Heads and Consistent Heads}
% \paragraph{Differences}
% Both Consistent Heads and Streaming Heads exploit distinctive attention patterns within the model. However, consistent Heads emphasize the overall similarity of attention across different sequences, while Streaming Heads prioritize whether initial and recent tokens cover most of the attention. In the process of identifying consistent Heads, CateKV considers the initial and recent tokens to be consistently important, classifying based on the remaining tokens. This is in contrast to the classification approach used for Streaming Heads. The attention map in Figure~\ref{fig:streaming} shows that there is no direct correlation between these two types. 

% \paragraph{Advantages}
% Consistent Heads incorporate additional priors from the pre-filling stage, enhancing their overall precision. In the right panel of Figure~\ref{fig:streaming}, we present attention weight recall curves for all attention heads, comparing the consistent and streaming patterns within the same example. The consistent pattern outperforms the streaming pattern and approaches the attention weight recall levels of the KV Retrieval method, such as Quest~\cite{quest}. As the proportion of streaming heads increases, attention recall declines noticeably. Therefore, the high precision of the consistent pattern can be combined with streaming-based methods to achieve a higher compression ratio.

\section{Experiments}
\subsection{Setup}
\paragraph{LLM and Benchmark}
We employed five widely used LLMs in long-context scenarios: LLaMA-3-8B-Instruct-1048K~\cite{llama3-8b}, Phi-3-Mini-128K~\cite{phi}, Llama-3.1-8B~\cite{llama3-1}, Yi-9B-200K~\cite{yi} and Qwen2.5-7B~\cite{qwen}. The performance of \shortname was assessed on three challenging benchmarks: RULER~\cite{ruler}, LongBench~\cite{longbench}, and Needle in a Haystack (NIAH)~\cite{niah}. We built a reference set for head identification of CateKV by emulating Variable Tracking task from RULER, which is very distinct from the test set. The experiments were carried out on a single NVIDIA A100-80G GPU.

\paragraph{Baselines}
% Considering their relevance to our method and code availability, 
We compare CateKV with eviction-based algorithms StreamingLLM~\cite{streamingllm}, SnapKV~\cite{snapkv}, PyramidKV~\cite{pyramidkv}, retrieval-based algorithms Quest~\cite{quest} and ShadowKV~\cite{shadowkv}, and the head-wise classification algorithm Duoattention~\cite{duoattention}.
In our experiments, all approaches maintained an exact pre-filling stage and utilized sparse attention during the decoding stage.  
We also do not perform memory optimization like ShadowKV in the accuracy comparison.
For fairness, when comparing with the KV eviction methods and head classification methods, we maintain the same KV cache size, while comparing with KV retrieval methods, we integrate them under adaptive heads to ensure an equivalent computational budget.
Given that Duoattention only provides attention patterns for Llama3 and Llama3.1 in the models used in our study, the comparison is restricted to these two models.
When baselines require a chunk size, we all set it to 8 to maintain consistency.
Further experimental details are in the Appendix~\ref{appendix-exp-detail}.




\subsection{Effectiveness Evaluation}
\label{acc_experiments}

% \input{Tables/longbench_2}
\subsubsection{Results on RULER}
% We evaluated the performance of our method on the RULER benchmark with a context length of 128K. 
In this experiment, we test 12 synthetic tasks under the context of 128K, with each task including 96 samples. To ensure a fair comparison with other baseline methods, we focus here on the task-aware setting.
% The Common Words Extraction (CWE) task is excluded because all the models performed poorly on this task, even when utilizing full attention. 
The results are shown in Table~\ref{tab:ruler-128k-eviction}. 
% Notably, we set the sparse budget at 2048 (1.56\%) on all sparse heads. 
% Specifically, we use a 40\% adaptive head ratio in \shortname and set the sparse budget of baselines to 40\% for a fair comparison.
Specifically, in \shortname, we set the adaptive head ratio $r$ to 0.4, the retention ratio $\eta$ to 1.0 and allocate a sparse budget for consistent heads at 2048 (1.56\%), retaining approximately 41\% of the KV cache size. 
Experimental results demonstrate that our method outperforms baselines and is comparable to full attention, despite evicting over half of the KV pairs.
Besides, \shortname exhibits outstanding performance in complex tasks such as multi-document QA and variable tracking while maintaining high accuracy in other tasks. Due to the space limitation, we place results of more context lengths and combination with other retrieval-based methods in Appendix~\ref{appendix-ruler}.

\begin{table}[t]
\centering
\caption{Performance (\%) of different models and various methods on LongBench. We present the average score of all 21 tasks. The ``Budget'' refers to the computational budget.
% \shortname outperforms other methods with 42\% KV Cache size.
Please refer to the Appendix~\ref{appendix-longbench} for detailed data.}

\setlength{\tabcolsep}{3.0pt} % 调整列间距
\small
\resizebox{0.48\textwidth}{!}{
\begin{tabular}{l|cc|c|c|c|c|c}
\toprule
\multicolumn{3}{c|}{Model}& \multicolumn{1}{c|}{LLama-3} & \multicolumn{1}{c|}{Phi-3} & \multicolumn{1}{c|}{Llama-3.1} &  \multicolumn{1}{c|}{Yi} &\multicolumn{1}{c}{Qwen2.5} \\
\midrule
Methods & Cache & Budget & \multicolumn{5}{c}{Average Performance}\\
\midrule
Full & 100\% & 100\% & 31.27 & 34.00& 33.68 & 33.02& 30.03\\
\quad +CateKV & 42\% & 42\% & 31.48 & 33.73 & 33.70 & 32.83 & 29.92\\
Quest & 100\% & 3\% & 30.90 & 33.11 & 33.20 & 32.82 & 29.30\\
\quad +CateKV & 42\% & 3\% & 31.03 & 33.29 & 33.38 & 32.79 & 29.33\\
ShadowKV & 100\% & 3\% & 30.77 & 32.53 & 33.03 & 32.41 & 28.57\\
\quad +CateKV & 42\% & 3\% & 30.94 & 32.45 & 32.96 & 32.21 & 28.51\\

\bottomrule
\end{tabular}}
\label{tab:longbench_2}
\vspace{-0.3cm}
\end{table}


\subsubsection{Results on LongBench}
LongBench~\cite{longbench} is a comprehensive long-context benchmark including 6 main categories and 21 diverse tasks.
Following ShadowKV~\cite{shadowkv}, we test samples with lengths exceeding 4096 tokens. In CateKV, we set the adaptive ratio $r$ and retention ratio $\eta$ to 0.4 and 1.0 respectively. The budget for consistent heads is set to 512.
% Due to the high accuracy maintained with a smaller budget on the Longbench dataset, in this experiment, we combine \shortname with the query-aware retrieval method Quest and ShadowKV in adaptive heads. We compare \shortname with all baseline methods while maintaining the same computational load. The sparse budget for all heads is set to 512, which corresponds to approximately 3\% of the average length of the Longbench benchmark. 
As shown in Table~\ref{tab:longbench_2}, CateKV enables a reduction in KV cache size to 42\% across five LLMs while maintaining an accuracy decline of no more than 0.3\% on average. 
% In certain tasks, such as Llama3-8B, a slight improvement in performance is even observable. 
Furthermore, when integrated with KV retrieval methods like Quest~\cite{quest} and ShadowKV~\cite{shadowkv}, CateKV facilitates a reduction in memory usage with only a minimal impact on accuracy, not exceeding a 0.2\% decrease, under the same computational budget.

\begin{figure}[t!]
\begin{center}
    \includegraphics[width=0.5\textwidth]{Figures/niah-llama3-catekv.pdf}
    \vspace{-0.3cm}
    \caption{Needle In A Haystack from 20K to 1M in Llama3}
    \label{fig:niah-llama3-catekv}
    \vspace{-0.5cm}
\end{center}
\end{figure}

\subsubsection{Results on Needle In A Haystack}
As shown in Figure~\ref{fig:niah-llama3-catekv}, on the Needle In A Haystack dataset, \shortname demonstrates its robust ability to accurately identify and extract relevant information from long contexts, ranging from 20K to 1M tokens, while reducing memory and computational cost by more than half. Additional NIAH test results on other models are available in the Appendix~\ref{appendix-niah}.

% \subsubsection{KV Cache Size-Accuracy Trade-Off}
% To provide a comprehensive demonstration of our method's performance, we show CateKV under various KV cache size on the RULER-128K benchmark. In CateKV, we maintained a retention ratio of 1.0 and adjusted the overall KV cache size by tuning the adaptive head ratio. Figure~\ref{fig:ours_duo} presents the results for Llama3-8B and Llama3.1-8B. CateKV consistently outperformed Duoattention and maintained robust performance relative to full attention, affirming its effectiveness and adaptability. Notably, as shown in Tables~\ref{tab:ruler-128k-eviction} and~\ref{tab:longbench_2}, all models demonstrated performance comparable to full attention with an adaptive head ratio of 0.4. For specific models, this ratio could be reduced further; for instance, Llama3-8B maintained comparable performance to full attention at a ratio of 0.3.

\begin{table}[t]
\centering
\caption{Performance of \shortname combined with Duoattention and MInference on the RULER benchmark. The results are tested on the Llama-3-8B-1M model.}
\resizebox{0.48\textwidth}{!}{
\begin{tabular}{l|cccccc|c}
\toprule
Methods & 8K& 16K& 32K & 64K & 128K& 256K& Avg.\\ 
\midrule
Duoattention & 78.32  & 77.93 & 69.77 & 66.27 & 60.37 & 58.07 & 68.46 \\
Duoattention w/ CateKV  & 89.93  & 91.08 & 89.96 & 86.44 & 84.57 & 81.59 & 87.26 \\
\midrule
MInference & 91.33 & 92.28 & 89.66 & 84.97 & 84.57 & 81.10 & 87.32 \\
MInference w/ CateKV  & 91.32  & 92.13 & 89.71 & 85.90 & 85.14 & 82.27 & 87.74 \\
	
\bottomrule
\end{tabular}}
\label{tab:minference_quest}
\vspace{-0.3cm}
\end{table}


\subsubsection{CateKV vs. Duoattention}
% \subsubsection{Comparison and Combination with Duoattention}
Duoattention~\cite{duoattention} classifies attention heads into retrieval heads and streaming heads, similar to our approach. 
For a comprehensive comparison, we evaluate \shortname against DuoAttention across varying ratios of adaptive/retrieval heads, ranging from 0.1 to 1.0.
Figure~\ref{fig:ours_duo} shows that Duoattention's accuracy notably declines when the proportion of retrieval heads falls below 0.3. 
In contrast, \shortname sustains performance close to full attention even when the ratio of adaptive heads decreases to 0.1.
This discrepancy arises from the fundamental difference in attention patterns between our consistent head and their streaming head.
The consistent head in CateKV captures similarity across different queries, whereas the streaming head primarily focuses on the coverage of initial and recent tokens.
This essential difference enables our method to complement DuoAttention, mitigating its performance degradation at low full attention head ratios.
To validate this, we conducted experiments presented in Table~\ref{tab:minference_quest}. Specifically, we set the retrieval head ratio in DuoAttention to 0.3, with the remaining 0.7 as streaming heads.
In contrast, ``Duoattention w/ CateKV'' replaces part of the streaming heads with consistent heads, resulting in a 0.2 consistent head ratio and a 0.5 streaming head ratio.
The results demonstrate that incorporating \shortname significantly improves DuoAttention’s performance at low full attention head ratios, further validating the effectiveness of our consistency pattern.

\begin{figure*}[t]
   \centering
   \subfigure{
    \includegraphics[width=0.45\linewidth]{Figures/efficiency_memory.pdf}
        \label{fig:efficiency_memory}
   }
   \hfill
    \subfigure{
    \includegraphics[width=0.45\linewidth]{Figures/efficiency_latency.pdf}
        \label{fig:efficiency_latency}
    }
    % \vspace{-0.3cm}
    \caption{Comparison of decoding efficiency between \shortname and full attention. As the context length increases, both the memory usage (a) and decoding latency (b) increase linearly, but CateKV exhibits a smaller slope compared to full attention. OOM indicates that the GPU memory limit (80G) is exceeded; corresponding data is obtained by extrapolation. 
    % Our method achieves up to 2.45$\times$ memory reduction for MHA and 1.65$\times$ for GQA models, along with up to 2.13$\times$ latency reduction for MHA and 1.5$\times$ for GQA models.
    }
    \vspace{-0.5cm}
    \label{fig:efficiency_decoding}
\end{figure*}

\subsubsection{Efficient Pre-filling Methods Integration}
\label{combination}
% \paragraph{Compatibility with KV Retrieval Methods}
% Since CateKV retains most KV pairs in adaptive heads, it can be seamlessly integrated with KV retrieval methods. We have combined CateKV with Quest~\cite{quest} and ShadowKV~\cite{shadowkv}, and the results are displayed in Tables~\ref{tab:longbench_2} and~\ref{tab:ruler-128k-appendix}. These results indicate that CateKV enables Quest and ShadowKV to reduce the KV cache size to 42\% while maintaining performance. In other words, the KV retrieval methods help CateKV reduce overall computational cost while preserving robust results in most tasks.
% \paragraph{Compatibility with KV Retrieval Methods}
% Since our method is orthogonal to other KV Retrieval approaches, it can be integrated with advanced KV Retrieval methods. 
% For compatibility with Quest~\cite{quest}, we implemented Quest's sparse attention pattern on adaptive heads, applying the same sparse budget as used for consistent heads. 
% Similarly, we integrated ShadowKV~\cite{shadowkv} with \shortname by applying ShadowKV to the adaptive heads.
% The experimental results are presented in~\jc{Table~\ref{tab:longbench}.}
% % The experimental results are summarized in Table~\ref{tab:minference_quest}.
% Compared with the standalone \shortname, integrating \shortname with Quest or ShadowKV results in slight performance degradation but significantly reduces computational overhead. In other words, \shortname enables Quest and ShadowKV to achieve substantial memory savings.

\shortname is designed to accelerate the decoding stage and can naturally integrate with prefilling acceleration methods.
% To evaluate the compatibility of our method with such approaches, we combined \shortname with the state-of-the-art prefilling acceleration method, MInference~\cite{minference}. 
We also combined \shortname with the efficient prefilling acceleration method, MInference~\cite{minference}, and tested it on RULER with contexts ranging from 8K to 256K.
% Specifically, we directly replaced our prefilling stage with MInference.
As shown in Table~\ref{tab:minference_quest}, combining \shortname with MInference does not lead to performance degradation and even a slight performance improvement has been observed.
This shows that sequential consistency within the heads is unaffected by the pre-filling inference patterns, highlighting the robustness of our method.
Moreover, the additional acceleration of the prefilling stage further improves the overall inference speed.


% \paragraph{Integration with Duoattention}
% Our method and DuoAttention categorize attention heads from different perspectives, leading to significant differences in their classifications. As illustrated in Figure~\ref{fig:ours_duo}, DuoAttention demonstrates robust performance when the retrieval head ratio is at least 50\%; however, its effectiveness markedly declines as the retrieval head ratio decreases. We have found that the consistent heads identified by our approach can further enhance DuoAttention's compression capabilities. Detailed experimental results are presented in Table~\ref{tab:minference_quest}.

\subsection{Efficiency Evaluation}
% \begin{table*}[t]
% \centering
% \caption{Efficiency comparison between CateKV and Full Attention on a single A100 GPU. In the single-sample inference task, we utilized texts of lengths 500K, 180K, and 650K as inputs for the following three models. For the batch-sample inference task, the sample length was set at 40K, and the maximum feasible batch size was used for each method evaluated.}
% \vspace{-0.2cm}
% \resizebox{0.95\textwidth}{!}{
% \begin{tabular}{l|cc|P{2.2cm}P{2.2cm}|P{2.2cm}P{2.2cm}}
% \toprule
%  & \multicolumn{2}{c|}{Accuracy (Avg)} & \multicolumn{2}{c|}{Single-sample inference} & \multicolumn{2}{c}{Batch-sample inference}\\
% \midrule
%  Model & RULER-128K & Longbench &Memory & Latency & Batchsize  & Throughout \\ 
% \midrule
% \textit{Llama-3-8B-1M} (8 KV heads)  & 84.10 (0.00) & 31.27 (0.00) & 77.72 ($1.00\times$) & 54.28 ($1.00\times$) & 12 ($1.00\times$) & 229.37 ($1.00\times$)\\
% CateKV ($r=0.4,\eta=1.0$)  & 84.61 (+0.51) & 31.48 (+0.21) & 41.02 ($1.89\times$) & 32.75 ($1.66\times$) & 28 ($2.33\times$) & 511.46 ($2.23\times$) \\
% CateKV ($r=0.3,\eta=0.7$) & 83.85 (-0.25) & 31.21 (-0.06) & 28.59 ($2.72\times$)& 24.86 ($2.18\times$)& & \\
% \midrule
% \textit{Phi-3-Mini-128K} (32 KV heads) & 72.06 (0.00)& 34.00 (0.00)& 75.11 ($1.00\times$) & 55.51 ($1.00\times$) & 4 ($1.00\times$) & 76.73 ($1.00\times$)\\
% CateKV ($r=0.4,\eta=1.0$) & 71.69 (-0.37)& 33.73 (-0.27) & 35.58 ($2.11\times$) & 31.06 ($1.79\times$) & & \\

% CateKV ($r=0.3,\eta=1.0$)  & 71.48 (-0.58) & 33.66 (-0.33) & 29.02 ($2.59\times$)& 27.81 ($2.00\times$) & & \\

% \midrule
% \textit{Yi-9B-200K} (4 KV heads) & 64.52 (0.00)& 33.02 (0.00)& 77.65 ($1.00\times$) & 56.30 ($1.00\times$) & 4 ($1.00\times$) & 76.73 ($1.00\times$)\\
% CateKV ($r=0.4,\eta=1.0$)  & 65.76 (+1.24) & 32.83 (-0.19) & 41.62 ($1.87\times$)& 36.05 ($1.56\times$)& & \\
% CateKV ($r=0.25,\eta=0.8$) & & & & & & \\

	
% \bottomrule
% \end{tabular}}
% \label{tab:efficiency}
% \vspace{-0.3cm}
% \end{table*}

\begin{table*}[t]
\centering
\caption{Efficiency comparison between CateKV and Full Attention on a single A100 GPU. In the single-sample inference task, we utilized texts of lengths 500K, 180K, and 650K as inputs for the following three models. For the batch-sample inference task, the sample length was set at 40K, and the maximum feasible batch size was used for each method evaluated.}
\resizebox{0.95\textwidth}{!}{
\begin{tabular}{l|cc|P{2.2cm}P{2.2cm}|P{2.2cm}P{2.2cm}}
\toprule
 & \multicolumn{2}{c|}{Accuracy (Avg)} & \multicolumn{2}{c|}{Single-sample inference} & \multicolumn{2}{c}{Batch-sample inference}\\
\midrule
 Model & RULER-128K & Longbench &Memory & Latency & Batchsize  & Throughput \\ 
\midrule
\textit{Llama-3-8B-1M} (8 KV heads)  & 84.10 \color{gray}{(0.00)} & 31.27 \color{gray}{(0.00)} & 77.72 \color{gray}{($1.00\times$)} & 54.28 \color{gray}{($1.00\times$)} & 12 \color{gray}{($1.00\times$)} & 229.37 \color{gray}{($1.00\times$)}\\
CateKV ($r=0.4,\eta=1.0$)  & 84.61 \color{gray}{(+0.51)} & 31.48 \color{gray}{(+0.21)} & 41.02 \color{gray}{($1.89\times$)} & 32.75 \color{gray}{($1.66\times$)} & 28 \color{gray}{($2.33\times$)} & 511.46 \color{gray}{($2.23\times$)} \\
CateKV ($r=0.3,\eta=0.7$) & 83.85 \color{gray}{(-0.25)} & 31.21 \color{gray}{(-0.06)} & 28.59 \color{gray}{($2.72\times$)}& 24.86 \color{gray}{($2.18\times$)}& 52 \color{gray}{($4.33\times$)}& 909.69
\color{gray}{($3.96\times$)} \\
\midrule
\textit{Phi-3-Mini-128K} (32 KV heads) & 72.06 \color{gray}{(0.00)}& 34.00 \color{gray}{(0.00)}& 75.11 \color{gray}{($1.00\times$)} & 55.51 \color{gray}{($1.00\times$)} & 4 \color{gray}{($1.00\times$)} & 78.53 \color{gray}{($1.00\times$)}\\
CateKV ($r=0.4,\eta=1.0$) & 71.69 \color{gray}{(-0.37)}& 33.73 \color{gray}{(-0.27)} & 35.58 \color{gray}{($2.11\times$)} & 31.06 \color{gray}{($1.79\times$)} & 10 \color{gray}{($2.50\times$)} & 167.11 \color{gray}{($2.13\times$)} \\

CateKV ($r=0.3,\eta=1.0$)  & 71.48 \color{gray}{(-0.58)} & 33.66 \color{gray}{(-0.33)} & 29.02 \color{gray}{($2.59\times$)}& 27.81 \color{gray}{($2.00\times$)} & 14 \color{gray}{($3.50\times$)} & 221.29 \color{gray}{($2.82\times$)} \\

\midrule
\textit{Yi-9B-200K} (4 KV heads) & 64.52 \color{gray}{(0.00)}& 33.02 \color{gray}{(0.00)}& 77.65 \color{gray}{($1.00\times$)} & 56.30 \color{gray}{($1.00\times$)} & 16 \color{gray}{($1.00\times$)} & 292.55 \color{gray}{($1.00\times$)}\\
CateKV ($r=0.4,\eta=1.0$)  & 65.76 \color{gray}{(+1.24)} & 32.83 \color{gray}{(-0.19)} & 41.62 \color{gray}{($1.87\times$)}& 36.05 \color{gray}{($1.56\times$)}& 38 \color{gray}{($2.38\times$)} & 659.02 \color{gray}{($2.25\times$)}\\
CateKV ($r=0.3,\eta=0.8$) & 65.75 \color{gray}{(+1.23)} & 32.78 \color{gray}{(-0.24)} & 31.58 \color{gray}{($2.46\times$)}& 29.96 \color{gray}{($1.88\times$)}& 60 \color{gray}{($3.75\times$)} & 980.86 \color{gray}{($3.35\times$)}\\

\bottomrule
\end{tabular}}
\label{tab:efficiency}
\vspace{-0.3cm}
\end{table*}


% We verify the efficiency of \shortname by evaluating the decoding latency and memory usage. 
% All the experiments are conducted on NVIDIA A100 GPUs.
% Overall, \shortname can achieve a 2.14$\times$reduction in memory usage and a 1.81$\times$acceleration during the decoding stage.
To validate the efficiency of CateKV, we tested it under both single-sample and batch-sample inference scenarios.

\subsubsection{Single-sample Inference}
To evaluate the efficiency of CateKV during single-sample inference, we selected three models with different numbers of KV heads: Llama-3, Phi-3, and Yi. We compared CateKV to full attention on these models by measuring memory usage and decoding latency at the same input length. 
We observed that memory and latency reductions increase as the context length grows (detailed in the Figure~\ref{fig:efficiency_decoding}).
To demonstrate optimal performance, we selected the maximum input length that can be handled by full attention on a single A100 GPU. As shown in Table~\ref{tab:efficiency}, under the generic settings of $r=0.4$ and $\eta=1.0$, the Phi-3 model achieved reductions of $2.11\times$ in memory and $1.79\times$ in latency by using CateKV. By balancing efficiency and accuracy, CateKV further reduced memory usage by $2.72\times$ and decoding latency by $2.18\times$ on Llama-3, with accuracy decline on RULER-128K and Longbench tasks under 0.25\%.


\subsubsection{Batch-sample Inference}
Batch-sample inference is a more common scenario in real-world applications. By evicting a portion of the KV pairs, CateKV supports larger batch sizes, thereby increasing throughput. To evaluate CateKV's efficiency in batch-sample scenarios, we set each sample length to 40K and input them in batches, comparing throughput at the maximum batch size for both full attention and CateKV. The comparative results across three models are presented in Table~\ref{tab:efficiency}. With the generic settings, gains of up to $2.38\times$ in batch size and $2.25\times$ in throughput were achieved. Additionally, through further KV cache compression while maintaining accuracy, both batch size and throughput in the Llama-3 model were boosted to $4.33\times$ and $3.96\times$, respectively.

\begin{figure}[t!]
\begin{center}
    \includegraphics[width=0.48\textwidth]{Figures/ours_duo_dark.pdf}
    \vspace{-0.2cm}
    \caption{Comparison of accuracy between our method and DuoAttention across different full attention head ratios.
    % Our method maintains performance comparable to full attention at lower full attention head ratios.
    }
    \vspace{-0.5cm}
    \label{fig:ours_duo}
\end{center}
\end{figure}

% Figure~\ref{fig:efficiency_decoding} illustrates the variations in memory consumption and decoding latency of \shortname across different context lengths. 
% In comparison to full attention, our approach demonstrates a substantial reduction in both decoding latency and memory utilization. Furthermore, these metrics exhibit a progressive decline as the context length increases, highlighting the scalability and efficiency of \shortname.

% \subsubsection{Ratio of Adaptive Heads}
% The adaptive head ratio $r$ is a pivotal control hyper-parameter that balances model accuracy and inference speed. 
% Specifically, higher values of 
% $r$ yield improved performance but greater decoding latency.
% Notably, the algorithm exhibits two boundary conditions: it degenerates to a KV cache eviction algorithm when $r$ reaches 0 and transforms into a full attention mechanism implementation when $r$ attains 1.
% To systematically analyze the impact of $r$, we evaluate the decoding latency and memory consumption of our method across a range of $r$ values. 
% (Corresponding performance variations can be found in~\ref{ablation:r-performance}.)
% As evidenced by the results in Figure~\ref{fig:ablation_ratio}, both decoding latency and memory usage exhibit a linear reduction as $r$ decreases.
% Dynamically adjusting $r$ in response to specific task requirements enables optimal balancing between model performance and computational efficiency, facilitating adaptive optimization tailored to diverse application scenarios.
% Under our settings in Table~\ref{tab:ruler-128k-eviction}, \shortname achieves maximum acceleration: 2.14$\times$ memory reduction and 1.81$\times$ latency reduction.



% \paragraph{Memory Usage}
% pass
% \begin{figure*}[t]
%     \centering
%     \includegraphics[width=\linewidth]{Figures/ablation_budget.pdf}
%     \caption{Comparison results between the models with full cache, \shortname, and Quest.}
%     \label{fig:ablation_budget}
% \end{figure*}

% \begin{figure}[t]
%    \centering
%    \subfigure{
%     \includegraphics[width=0.46\linewidth]{Figures/llama3-128k-sparse.pdf}
%         \label{fig:budget_l3}
%    }
%    \hfill
%     \subfigure{
%     \includegraphics[width=0.46\linewidth]{Figures/llama3-128k-sparse.pdf}
%         \label{fig:budget_l3_2}
%     }
%     \caption{Impact of sparse budget on (a) Llama-3-8B-1M and (b) GLM-4-9B-1M.}
%     \vspace{-0.3cm}
%     \label{fig:ablation_budget}
% \end{figure}

% \begin{figure}[t]
%    \centering
%    \subfigure{
%     \includegraphics[width=0.46\linewidth]{Figures/ablation_needupdate.png}
%         \label{fig:chunk_acc}
%    }
%    \hfill
%     \subfigure{
%     \includegraphics[width=0.46\linewidth]{Figures/ablation_needupdate.png}
%         \label{fig:chunk_hitrate}
%     }
%     \caption{(a) Impact of chunk size on batch size and accuracy. (b) Minimal effect of chunk size on hit rate.}
%     \vspace{-0.3cm}
%     \label{fig:ablation_ratio}
% \end{figure}

\begin{figure*}[t]
\begin{center}
    \includegraphics[width=0.95\textwidth]{Figures/Ablation.pdf}
    \vspace{-0.2cm}
    \caption{(a) The effect of adaptive head ratio on memory and decoding latency is approximately linear. (b) Impact of retention ratio on accuracy in RULER-128K and decoding latency for 500K length input. (c) Minimal effect of sparse budget in consistent heads on accuracy.}
    \vspace{-0.5cm}
    \label{fig:ablation}
\end{center}
\end{figure*}


\subsection{Ablation Study} 
The ablation studies focus on three key aspects of \shortname: (1) ratio of adaptive heads, (2) retention ratio in the adaptive head, and (3) sparse KV cache budget in the consistent head.
% More ablation studies about different identification methods are in the appendix~\ref{different_methods_ablation}.
All experiments are conducted on the RULER-128K datasets using the Llama-3-8B-1M model.

\subsubsection{Ratio of Adaptive Heads}
\label{ablation:r-performance}
The ratio of adaptive heads $r$ is a crucial hyperparameter balancing accuracy and inference speed. As shown in Figure~\ref{fig:ours_duo}, the relationship between $r$ and model performance does not exhibit a perfect inverse correlation. 
Rather, the performance remains relatively stable with decreasing $r$ until reaching a critical threshold ($\approx$0.2), beyond which significant degradation occurs due to excessive reliance on critical tokens during the prefilling stage. In terms of efficiency, as shown in Figure~\ref{fig:ablation}(a), both memory usage and decoding latency decrease almost linearly with decreasing $r$.

% Notably, optimal performance is typically achieved around 
% $r=0.4$, indicating that moderate KV cache sparsification can be beneficial for enhancing model performance while improving efficiency.

\subsubsection{Retention Ratio in Adaptive Head}
In adaptive heads, there is still a subset of tokens that are always not important.
This enables a reduction in the proportion of the KV cache retained within adaptive heads. Figure~\ref{fig:ablation}(b) illustrates the impact of changes in the retention ratio $\eta$ on accuracy and latency.
When $\eta$ exceeds a certain threshold ($\approx$0.6), the model maintains or even surpasses the performance achieved with the full KV cache. 
However, dropping below this threshold results in a significant performance decline, indicating that adaptive heads heavily depend on the majority of the KV cache. Additionally, a decrease in $\eta$ leads to a linear reduction in latency.
In practical applications, it is essential to adjust both $r$ and $\eta$ to balance accuracy and memory consumption.

\subsubsection{Sparse Budget in Consistent Head}

As illustrated in Figure~\ref{fig:ablation}(c), \shortname demonstrates strong robustness across different sparse budgets. Under the settings of $r=0.4$ and $\eta=1.0$, even when the sparse budget in the consistent head is reduced to approximately 0.78\% (1024), CateKV still maintains comparable performance to full attention in terms of average accuracy. 
 This indicates that, during inference, the consistent head only requires a minimal portion of the cache to perform its function, while the acquisition of global information relies primarily on adaptive heads, which retain the majority of KV pairs during the decoding stage. Since the sparse budget for the consistent head is a small fraction of the total cache, its impact on memory usage and inference latency is negligible.
% As illustrated in Figure~\ref{fig:ablation_budget}, \shortname presents strong robustness across different sparse budgets, consistently outperforming SOTA baselines.
% Our method can achieve comparable or even superior performance with full attention under most settings.
% Experimental observations reveal that all sparse methods are characterized by a critical sparse budget threshold, beyond which substantial performance deterioration occurs.
% It can be observed that the proposed method achieves a lower critical sparse budget threshold, which can be attributed to its adaptive heads that preserve full attention in the decoding stage.


% In this ablation study, we maintain the sparse budget size consistent with that described in Section~\ref{acc_experiments}.

\section{Conclusion}
We propose CateKV, a novel hybrid KV cache method that leverages sequential consistency to improve LLM inference efficiency in long-context tasks. By using a coefficient-of-variation-based algorithm, CateKV classifies attention heads into consistent and adaptive types. It selectively retains critical KV pairs in consistent heads and most pairs in adaptive heads, reducing memory usage and decoding latency while maintaining performance. Additionally, it can be easily integrated with other acceleration methods for further enhancement. Extensive evaluations demonstrate that CateKV achieves significant efficiency gains, including up to $2.72\times$ reduction in memory usage, $2.18\times$ acceleration in decoding, and a $3.96\times$ throughput increase in batch scenarios.

\section*{Acknowledgments}
This work is supported by National Natural Science Foundation of China (No.62306178), STCSM (No.22DZ2229005), 111 plan (No. BP0719010).
\section*{Impact Statement}
The primary objective of our research is to develop inference acceleration methods for large language models, aiming to advance their efficiency and scalability. We explicitly state that any depictions of violence in the datasets are wholly fictional and are utilized exclusively for academic research purposes. This work does not reflect, support, or justify any real-world violent actions. Additionally, this research was carried out independently, with no external funding or conflicts of interest influencing its outcomes. The study strictly follows ethical guidelines, taking into account essential factors such as discrimination, bias, fairness, privacy, security, and legal adherence, while maintaining the utmost integrity in the research process.

\bibliography{ref}
\bibliographystyle{icml2025}


% APPENDIX

\newpage
\appendix
\onecolumn
\section{Further Observations}
\label{appendix-consistency}

We expanded our exploration of sequential consistency to a wider selection of models, primarily focusing on the popular models mentioned in the main text: Llama-3-8B-1M~\cite{llama3-8b}, Llama-3.1-8B~\cite{llama3-1}, Phi-3-128K~\cite{phi}, and Yi-9B-200K~\cite{yi}. In Figure~\ref{fig:appendix-consistency}, we visualized the attention weights heatmap for these four models, illustrating the presence of both consistent and adaptive heads across various layers. This visualization supports the generality of our observations. Consistent with the setup described in the main text, the heatmap in the figure comprises attention weights associated with the last 20 query tokens during the pre-filling stage and all query tokens during the decoding stage, employing a causal mask. And the samples are randomly excerpted from the WikiText-2~\cite{wikitext} dataset.


\section{Experiment Details}
\subsection{Implementation Details of Experiments}
\label{appendix-exp-detail}
During the identification stage of CateKV, we employed an observation window and temporarily excluded initial tokens and recent tokens from the context window. We set $L_{obs}$ to 64, while $L_{init}$ and $L_{rec}$ were defined as 1/32 and 1/128 of the sparse budget, respectively. We constructed the reference dataset based on the Variable Tracking task from the RULER Benchmark, which comprises 100 samples, each with a length of 128K, distinct from the test set. The sparse budget was set at 2048. According to the performance on the reference dataset, we selected the most appropriate percentile threshold $k$ and scaling factor $\alpha$ for each model. For the percentile threshold $k$, Llama3 and Llama3.1 were set at 0.996 and 0.984, respectively, while other models were set at 0.99. For the scaling factor $\alpha$, Llama3.1, and Yi were set at 0.8, while other models were assigned a value of 1.0.

For the baseline methods, we configured the observation windows of SnapKV~\cite{snapkv} and PyramidKV~\cite{pyramidkv} to 32 and set the $\beta$ in PyramidKV to 20. For StreamingLLM~\cite{streamingllm}, the initial tokens were set to 128. Regarding Duoattention\cite{duoattention}, we conducted experiments using the attention patterns provided by their code available on GitHub.

\subsection{Additional Results on RULER}
\label{appendix-ruler}
\subsubsection{Performance of different context lengths on RULER}
We also conducted evaluations on various context lengths within the RULER benchmark. The Table ~\ref{tab:ruler8k-256k} presents the performance of CateKV on tasks with context lengths from 8K to 256K. CateKV is comparable to full attention in terms of all lengths and average results and even shows slight improvements in performance at certain lengths when $r=0.4$ and $\eta=1.0$.
\begin{table}[h]
\centering
\caption{Performance of different context lengths on RULER}
\begin{tabular}{l|cccccc|c}
\toprule
Methods & 8K& 16K& 32K & 64K & 128K& 256K& Avg.\\ 
\midrule
\textit{Llama-3-8B-1M} & 91.47 & 92.87 & 90.31 & 86.44 & 84.10 & 79.79 & 87.50 \\
CateKV  & 91.28 & 92.11 & 90.37 & 86.86 & 84.61 & 81.53 & 87.79 \\
\midrule
\textit{Phi-3-Mini-128K} & 92.02 & 91.42 & 91.24 & 87.89 & 72.06 & - & 86.93 \\
CateKV  & 92.68 & 92.54 & 92.04 & 88.78 & 71.69 & - & 87.55 \\
\midrule
\textit{Llama-3.1-8B} & 94.78 & 94.95 & 94.61 &  93.02 & 84.55 & - & 92.38\\
CateKV  & 94.66 & 94.68 & 94.58 & 93.03 & 84.66 & - & 92.32 \\
\midrule
\textit{Yi-9B-200K} & 87.54 & 82.33 & 72.06 & 69.19 & 64.52 & - & 75.13 \\
CateKV  & 87.17 & 81.68 & 71.85 & 69.25 & 65.76 & - &  75.14\\
	
\bottomrule
\end{tabular}
\label{tab:ruler8k-256k}
\end{table}

\subsubsection{Combine with KV Retrieval Methods}
We combine CateKV with Quest~\cite{quest} and ShadowKV~\cite{shadowkv} and compare their performance with the baseline under the same computational budget at a length of 128K. The results are shown in Table~\ref{tab:ruler-128k-appendix}. CateKV helps reduce the memory usage of Quest and ShadowKV and significantly outperforms KV Eviction methods in terms of accuracy with the same computational load.
\begin{table*}[t]
\centering
\caption{Performance (\%) of CateKV combined with KV retrieval methods. 'Budget' refers to the computational budget for sparse attention. CateKV can reduce memory. CateKV can help KV retrieval methods reduce memory usage to 41\% while maintaining accuracy. }

\setlength{\tabcolsep}{2pt} % 调整列间距
\small

\begin{tabular}{l|c|cccccccccccc|c}
\toprule
Methods & Budget & N-S1 & N-S2 & N-S3 & N-MK1 & N-MK2 & N-MK3 & FWE & N-MQ & N-MV & QA-1 & QA-2 & VT & Avg. \\ 
\midrule
\textit{Llama-3-8B-1M} & 100\% & 100.00 & 100.00 & 100.00 & 98.96 & 98.96 & 41.67 & 71.88 & 98.69 & 96.35 & 73.96 & 50.00 & 78.75 & 84.10 \\
SnapKV & 1.6\% & 100.00 & 100.00 & 14.58 & 98.96 & 96.88 & 0.00 & 61.11 & 98.44 & 96.88 & 68.75 & 48.96 & 79.38 & 72.00 \\
PyramidKV & 1.6\% & 100.00 & 100.00 & 10.42 & 98.96 & 96.88 & 0.00 & 56.60 & 98.18 & 95.58 & 70.83 & 48.96 & 80.42 & 71.40 \\
Quest & 1.6\% & 100.00 & 100.00 & 100.00 & 98.96 & 97.92 & 19.79 & 58.33 & 98.96 & 96.61 & 72.91 & 52.08 & 80.20 & 81.31 \\
CateKV+Quest & 1.6\% & 100.00 & 100.00 & 100.00 & 97.92 & 97.92 & 19.79 & 58.33 & 98.96 & 94.27 & 73.96 & 51.04 & 82.29 & 81.21 \\
ShadowKV & 1.6\% & 100.00 & 100.00 & 100.00 & 97.92 & 93.75 & 21.88 & 75.69 & 98.96 & 96.09 & 72.92 & 50.00 & 78.96 & 82.18 \\
CateKV+ShadowKV & 1.6\% & 100.00 & 100.00 & 100.00 & 97.92 & 89.58 & 21.88 & 75.35 & 98.70 & 94.79 & 71.88 & 51.04 & 81.25 & 81.86 \\
\midrule
\textit{Phi-3-Mini-128K} & 100\% & 96.88 & 90.63 & 95.83 & 83.33 & 65.63 & 37.50 & 87.15 & 72.14 & 66.67 & 63.54 & 39.58 & 65.83 & 72.06 \\
SnapKV & 1.6\% & 98.21 & 38.54 & 1.04 & 42.71 & 11.46 & 0.00 & 60.76 & 8.59 & 2.60 & 62.50 & 38.54 & 60.63 & 35.47 \\
PyramidKV & 1.6\% & 97.92 & 39.58 & 0.00 & 46.88 & 11.46 & 0.00 & 56.94 & 10.16 & 2.08 & 59.38 & 39.58 & 60.42 & 35.37 \\
Quest & 1.6\% & 96.88 & 92.71 & 96.88 & 80.21 & 57.29 & 20.83 & 57.29 & 69.53 & 63.02 & 64.58 & 39.58 & 63.75 & 66.88 \\
CateKV+Quest & 1.6\% & 96.88 & 92.71 & 96.88 & 80.21 & 57.29 & 18.75 & 55.56 & 68.75 & 64.32 & 64.58 & 39.58 & 62.50 & 66.50 \\
ShadowKV & 1.6\% & 95.83 & 88.54 & 90.63 & 80.21 & 54.17 & 21.88 & 77.43 & 63.28 & 51.28 & 62.50 & 38.54 & 63.75 & 65.71 \\
CateKV+ShadowKV & 1.6\% & 97.92 & 87.50 & 92.70 & 77.08 & 55.20 & 18.75 & 74.31 & 64.06 & 55.99 & 62.50 & 38.54 & 65.83 & 65.87 \\
\midrule
\textit{Llama-3.1-8B} & 100\% & 100.00 & 100.00 & 98.96 & 98.96 & 90.63 & 63.54 & 71.53 & 98.96 & 95.31 & 81.25 & 46.88 & 68.54 & 84.55 \\
SnapKV & 1.6\% & 100.00 & 100.00 & 41.67 & 98.96 & 79.17 & 0.00 & 59.72 & 97.14 & 91.67 & 81.25 & 44.79 & 62.92 & 71.44 \\
PyramidKV & 1.6\% & 100.00 & 100.00 & 33.33 & 98.96 & 81.25 & 1.04 & 56.25 & 95.83 & 93.48 & 81.25 & 45.83 & 65.00 & 71.02 \\
Quest & 1.6\% & 100.00 & 100.00 & 100.00 & 98.96 & 78.13 & 4.17 & 59.03 & 98.70 & 94.01 & 80.21 & 50.00 & 68.96 & 77.68 \\
CateKV+Quest & 1.6\% & 100.00 & 100.00 & 100.00 & 98.96 & 78.13 & 3.13 & 62.15 & 97.92 & 90.63 & 82.29 & 47.92 & 66.88 & 77.33 \\
ShadowKV & 1.6\% & 100.00 & 100.00 & 98.96 & 98.96 & 77.08 & 13.54 & 70.49 & 98.18 & 90.36 & 81.25 & 48.96 & 64.17 & 78.50 \\
CateKV+ShadowKV & 1.6\% & 100.00 & 100.00 & 100.00 & 98.96 & 73.96 & 13.54 & 65.97 & 98.44 & 90.89 & 80.21 & 50.00 & 67.08 & 78.25 \\
\midrule
\textit{Yi-9B-200K} & 100\% & 100.00 & 100.00 & 98.96 & 85.42 & 63.54 & 18.75 & 89.24 & 66.41 & 32.55 & 45.83 & 38.54 & 35.00 & 64.52 \\
SnapKV & 1.6\% & 100.00 &	93.75 &	5.21 &	80.21 &	6.25 &	0.00 	&75.69 &	54.95 &	18.23 &	40.63 &	36.46 &	52.92 &	47.02 \\
PyramidKV & 1.6\% & 100.00 & 94.79 & 4.17 & 79.17 & 5.21 & 0.00 & 85.76 & 55.73 & 15.63 & 41.67 & 33.33 & 51.45 & 47.24 \\
Quest & 1.6\% & 100.00 & 97.92 & 98.96 & 85.42 & 64.58 & 4.17 & 67.01 & 66.14 & 38.39 & 42.71 & 36.46 & 48.33 & 62.51 \\
CateKV+Quest & 1.6\% & 100.00 & 100.00 & 98.96 & 85.42 & 70.83 & 5.21 & 66.67 & 63.28 & 37.24 & 40.63 & 37.50 & 50.21 & 62.99 \\
ShadowKV & 1.6\% & 100.00 & 100.00 & 97.92 & 87.50 & 60.42 & 2.08 & 75.35 & 59.64 & 34.11 & 43.75 & 37.50 & 54.79 & 62.76 \\
CateKV+ShadowKV & 1.6\% & 100.00 & 100.00 & 97.92 & 84.38 & 58.33 & 2.08 & 82.29 & 59.38 & 33.85 & 43.75 & 36.46 & 47.29 & 62.14 \\
\bottomrule
\end{tabular}
\label{tab:ruler-128k-appendix}
\vspace{-0.3cm}
\end{table*}


\subsection{Full Longbench Results}
\label{appendix-longbench}
We present the complete experimental results for the Longbench in Table~\ref{tab:longbench-all}. We integrate CateKV with both the Full Attention and KV retrieval methods, Quest\cite{quest} and ShadowKV\cite{shadowkv}, and evaluate its performance on all 21 tasks in Longbench. The results showed that this integration did not lead to any significant drop in per-task accuracy, and the average accuracy even outperformed the original methods, despite retaining only 42\% of the KV cache size. For around half of the tasks, the combination of CateKV with the original methods leads to a slight improvement in performance.

\begin{table*}[t]
\centering
\caption{Full LongBench results with Llama-3-8B-1M.}
\label{tab:transposed-longbench}
\begin{tabular}{l|cc|cc|cc}
\toprule
Metrics & Full & CateKV  & Quest  & CateKV+Quest & ShadowKV & CateKV+ShadowKV  \\ 
\midrule
Average & 31.27 & \textbf{31.48} & 30.90 & \textbf{31.03} & 30.77 & \textbf{30.94} \\
\midrule
NarrativeQA & 18.61 & \textbf{18.64} & \textbf{19.54} & 17.93 & \textbf{18.43} & 17.74 \\
Qasper & 25.83 & \textbf{27.00} & \textbf{27.00} & 26.20 & 25.39 & \textbf{26.38} \\
MultiFieldQA-en & 48.06 & \textbf{48.17} & \textbf{45.80} & 45.72 & 45.59 & \textbf{46.63} \\
MultiFieldQA-zh & \textbf{33.76} & 33.68 & \textbf{34.23} & 33.37 & \textbf{34.23} & 33.65 \\
HotpotQA & 36.35 & \textbf{36.44} & 35.79 & \textbf{36.84} & \textbf{38.00} & 37.64 \\
2WikiMultihopQA & \textbf{25.17} & 24.61 & \textbf{25.48} & 24.07 & 24.92 & \textbf{25.98} \\
MuSiQue & \textbf{21.08} & 20.41 & \textbf{20.18} & 19.56 & \textbf{20.70} & 20.22 \\
DuReader & \textbf{30.98} & 28.62 & \textbf{29.23} & 27.36 & \textbf{29.82} & 27.91 \\
GovReport & 23.38 & \textbf{23.45} & \textbf{23.96} & 23.53 & 22.35 & \textbf{22.85} \\
QMSum & \textbf{25.45} & 24.74 & 24.59 & \textbf{24.66} & \textbf{24.67} & 24.29 \\
MultiNews & \textbf{22.63} & 21.42 & \textbf{23.30} & 21.10 & \textbf{23.53} & 21.10 \\
VCSUM & \textbf{14.19} & 13.90 & \textbf{14.21} & 13.93 & \textbf{13.86} & \textbf{13.86} \\
TREC & 39.00 & \textbf{41.83} & 39.11 & \textbf{39.49} & 37.69 & \textbf{39.49} \\
TriviaQA & 16.81 & \textbf{16.91} & 16.67 & \textbf{16.73} & \textbf{17.08} & 16.91 \\
SAMSum & 26.46 & \textbf{26.59} & \textbf{26.66} & 25.16 & \textbf{26.02} & 25.61 \\
LSHT & 31.58 & \textbf{32.00} & 24.75 & \textbf{33.25} & 29.68 & 29.96 \\
PassageCount & \textbf{1.00} & \textbf{1.00} & \textbf{1.00} & \textbf{1.00} & \textbf{1.00} & \textbf{1.00} \\
PassageRetrieval-en & \textbf{81.00} & 80.50 & 74.50 & \textbf{80.50} & 80.00 & \textbf{80.50} \\
PassageRetrieval-zh & \textbf{43.73} & 43.61 & \textbf{42.85} & 42.50 & \textbf{39.39} & 38.98 \\
LCC & 48.31 & \textbf{50.08} & \textbf51.30 & \textbf{51.45} & 49.06 & \textbf{50.69} \\
RepoBench-P & 43.21 & \textbf{47.58} & \textbf{48.73} & 48.27 & 43.84 & \textbf{48.45} \\
\bottomrule
\end{tabular}
\label{tab:longbench-all}
\end{table*}

\subsection{Additional Results in Needle In A Hystack}
\label{appendix-niah}

Figure~\ref{appendix-niah} displays the performance of the Llama3.1-8B, Phi-3-Mini-128K, and Yi-9B-200K models on the 'Needle In A Haystack' task. Compared to full attention, CateKV shows varied performance across different context windows and needle depths, but maintains overall comparable performance. This suggests that CateKV does not significantly affect the models' capacity to access and retrieve long-context semantic information.

\subsection{Additional Results on Larger Models}
We evaluated the performance of CateKV on larger models, setting the context length according to the maximum supported by each model—128k for Qwen2.5-32B and Yi-34B-200K, and 16k for Phi-4-14B. As shown in Table~\ref{tab:larger_models}, CateKV scales effectively, achieving near full-attention accuracy on the 30B and 14B models, outperforming baseline methods such as SnapKV and PyramidKV.

\begin{table*}[t]
\centering
\caption{Performance (\%) of CateKV on larger models}
\setlength{\tabcolsep}{3pt} % 调整列间距
\small

\begin{tabular}{l|c|cccccccccccc|c}
\toprule
Methods & Cache & N-S1 & N-S2 & N-S3 & N-MK1 & N-MK2 & N-MK3 & FWE & N-MQ & N-MV & QA-1 & QA-2 & VT & Avg. \\ 
\midrule
\textit{Qwen2.5-32B} & 100\% & 100.00 & 87.50 &	97.92 &	70.83 &	15.63 &	7.29 &	90.28 &	87.24 & 85.16 &	51.04 &	41.67 &	85.41 &	68.33 \\
SnapKV  & 41\% & 100.00 & 88.54 & 51.04 & 69.79 & 12.50 & 2.08 & 88.54 & 76.82 & 76.82 & 51.04 & 41.67 & 85.83 & 62.06 \\
PyramidKV  & 41\% & 100.00 & 87.50 & 46.88 & 66.67 & 8.33 &	1.04 &	84.02 &	66.93 &	67.71 &	48.96 &	41.67 &	84.79 &	58.71 \\
\rowcolor{cyan!10}
CateKV  & 41\% & 100.00 & 86.46 & 95.83 & 71.88 & 14.58 & 6.25 & 89.58 & 86.88 & 86.28 & 50.00 & 43.75 & 86.67 & 68.18 \\
\midrule
\textit{Yi-34B-200K} & 100\% & 100.00 & 100.00 & 100.00 & 92.71 & 70.83 & 47.92 & 86.11 & 97.14 & 92.45 & 68.75 & 47.92 & 88.05 & 82.66 \\
SnapKV & 41\% & 100.00 & 97.92 & 80.21 & 90.62 & 22.92 & 17.71 & 81.25 & 91.15 & 72.14 & 67.71 & 47.92 & 86.25 & 71.32 \\
PyramidKV & 41\% & 100.00 & 100.00 & 68.75 & 91.67 & 26.04 & 12.50 & 82.99 & 91.15 & 79.43 & 69.79 & 47.92 & 86.46 & 71.39 \\
\rowcolor{cyan!10}
CateKV & 41\% & 100.00 & 100.00 & 100.00 & 92.71 & 73.96 & 47.92 & 85.12 & 97.14 & 91.15 & 67.71 & 46.88 & 87.08 & 82.47 \\
\midrule
\textit{Phi-4-14B} & 100\% & 100.00 & 97.92 & 100.00 & 100.00 & 97.92 & 100.00 & 98.96 & 98.96 & 99.22 & 80.21 & 67.71 & 100.00 & 95.08 \\
SnapKV & 43\% & 100.00 & 100.00 & 7.29 & 100.00 & 93.75 & 3.13 & 99.31 & 97.66 & 99.22 & 82.29 & 66.67 & 100.00 & 79.11 \\
PyramidKV & 43\% & 100.00 & 100.00 & 3.13 & 100.00 & 94.79 & 5.21 & 98.96 & 98.44 & 98.44 & 80.21 & 67.71 & 100.00 & 78.91 \\
\rowcolor{cyan!10}
CateKV & 43\% & 100.00 & 98.96 & 100.00 & 97.92 & 98.96 & 100.00 & 99.31 & 98.44 & 99.48 & 78.13 & 67.71 & 99.79 & 94.89 \\
\bottomrule
\end{tabular}
\label{tab:larger_models}
\vspace{-0.3cm}
\end{table*}

% \begin{figure}[t]
%    \centering
%    \subfigure{
%     \includegraphics[width=0.48\linewidth]{Figures/efficiency_memory.pdf}
%         \label{fig:efficiency_memory}
%    }
%    \hfill
%     \subfigure{
%     \includegraphics[width=0.48\linewidth]{Figures/efficiency_latency.pdf}
%         \label{fig:efficiency_latency}
%     }
%     % \vspace{-0.3cm}
%     \caption{Comparison of decoding efficiency between \shortname and full attention. We compare the memory usage (a) and decoding latency (b) across different context lengths. OOM indicates that the GPU memory limit (80G) is exceeded; corresponding data is obtained by extrapolation. (can be found in Appendix)
%     % Our method achieves up to 2.45$\times$ memory reduction for MHA and 1.65$\times$ for GQA models, along with up to 2.13$\times$ latency reduction for MHA and 1.5$\times$ for GQA models.
%     }
%     \vspace{-0.5cm}
%     \label{fig:efficiency_decoding}
% \end{figure}

% \begin{figure}[t]
%    \centering
%    \subfigure{
%     \includegraphics[width=0.46\linewidth]{Figures/niah-llama3.1-full.pdf}
%         \label{fig:niah-llama3.1-full}
%    }
%    \hfill
%     \subfigure{
%     \includegraphics[width=0.46\linewidth]{Figures/niah-llama3.1-catekv.pdf}
%         \label{fig:niah-llama3.1-catekv}
%     }
%     \caption{Llama-3.1-8B-Instruct}
%     \vspace{-0.3cm}
%     \label{fig:niah-llama3.1}
% \end{figure}

% \subsection{Efficiency Across Different Context Lengths}
% \label{appendix-context-length}
% We compared the decoding efficiency of CateKV and full attention at different context lengths, as shown in Figure~\ref{fig:efficiency_decoding}. As the context length increases, both the model's latency and memory usage increase linearly, but CateKV exhibits a smaller slope compared to full attention. The longer the input context, the greater the reductions in memory and latency, ultimately approaching an upper bound, which theoretically would be $1/(r*\eta)$.

% \input{Tables/longbench_new}

% \input{Tables/3method}
% \subsubsection{Different identification scores}
% \label{different_methods_ablation}
% In Section~\ref{identification_method}, we introduce three scores for identifying consistent heads. We applied methods based on these scores to the same reference dataset to generate head classification masks and evaluated them on RULER-128K. Table~\ref{tab:3method} presents the results of the three methods under consistency ratios ranging from 0.3 to 0.5. While all three methods achieve performance comparable to full attention, the CV score-based method demonstrates superior performance.

% \begin{figure}[t]
%    \centering
%    \subfigure[Llama-3-8B-1M]{
%     \includegraphics[width=0.90\linewidth]{Figures/llama3_consistency.pdf}
%     \label{fig:llama3-consistency}
%    }
%    \vspace{0.01cm}
%    \subfigure[Llama-3.1-8B]{
%     \includegraphics[width=0.90\linewidth]{Figures/llama3.1_consistency.pdf}
%     \label{fig:llama3.1_consistency}
%    }
%    \vspace{0.01cm}  % Space between rows
%    \subfigure[Phi-3-Mini-128K]{
%     \includegraphics[width=0.90\linewidth]{Figures/Phi_consistency.pdf}
%     \label{fig:Phi_consistency}
%    }
%    \vspace{0.01cm}
%    \subfigure[Yi-9B-200K]{
%     \includegraphics[width=0.9\linewidth]{Figures/Yi_consistency.pdf}
%     \label{fig:Yi_consistency}
%    }
%     \caption{Sequential consistency in Llama-3-8B-1M~\cite{llama3-8b}, Llama-3.1-8B~\cite{llama3-1}, Phi-3-Mini-128K~\cite{phi} and Yi-9B-200K~\cite{yi}}
%     \vspace{-0.01cm}
%     \label{fig:appendix-consistency}
% \end{figure}



\begin{figure}[t]
   \centering
   \subfigure[Llama-3-8B-1M]{
    \includegraphics[width=0.90\linewidth]{Figures/llama3_consistency.pdf}
    \label{fig:llama3-consistency}
   }
   \vspace{0.01cm}
   \subfigure[Llama-3.1-8B]{
    \includegraphics[width=0.90\linewidth]{Figures/llama3.1_consistency.pdf}
    \label{fig:llama3.1_consistency}
   }
   \vspace{0.01cm}  % Space between rows
   \subfigure[Phi-3-Mini-128K]{
    \includegraphics[width=0.90\linewidth]{Figures/Phi_consistency.pdf}
    \label{fig:Phi_consistency}
   }
   \vspace{0.01cm}
   \subfigure[Yi-9B-200K]{
    \includegraphics[width=0.9\linewidth]{Figures/Yi_consistency.pdf}
    \label{fig:Yi_consistency}
   }
    \caption{Sequential consistency in Llama-3-8B-1M, Llama-3.1-8B, Phi-3-Mini-128K and Yi-9B-200K}
    \vspace{-0.01cm}
    \label{fig:appendix-consistency}
\end{figure}

\begin{figure}[t]
   \centering
   \subfigure[Llama-3.1-8B-Instruct]{
    \includegraphics[width=0.45\linewidth]{Figures/niah-llama3.1-full.pdf}
    \label{fig:niah-llama3.1-full}
   }
   \hfill
   \subfigure[Llama-3.1-8B-Instruct w/ CateKV]{
    \includegraphics[width=0.45\linewidth]{Figures/niah-llama3.1-catekv.pdf}
    \label{fig:niah-llama3.1-catekv}
   }
   \vspace{0.2cm}  % Space between rows
   \subfigure[Phi-3-mini-128k-instruct]{
    \includegraphics[width=0.45\linewidth]{Figures/niah-phi-full-2.pdf}
    \label{fig:niah-phi-full}
   }
   \hfill
   \subfigure[Phi-3-mini-128k-instruct w/ CateKV]{
    \includegraphics[width=0.45\linewidth]{Figures/niah-phi-catekv-2.pdf}
    \label{fig:niah-phi-catekv}
   }
   \subfigure[Yi-9B-200k]{
    \includegraphics[width=0.45\linewidth]{Figures/niah-yi-full.pdf}
    \label{fig:niah-yi-full}
   }
   \hfill
   \subfigure[Yi-9B-200k w/ CateKV]{
    \includegraphics[width=0.45\linewidth]{Figures/niah-yi-catekv.pdf}
    \label{fig:niah-yi-catekv}
   }
    \caption{NIAH Results on Llama-3.1-8B-Instruct\cite{llama3-1},  Phi-3-mini-128k-instruct~\cite{phi} and Yi-9B-200k~\cite{yi}}
    \vspace{-0.3cm}
    \label{fig:appendix-niah}
\end{figure}

% \begin{figure}[t]
% \begin{center}
%     \includegraphics[width=0.98\textwidth]{Figures/tasks_lengths.pdf}
%     \caption{Consistent pattern across different tasks and context lengths for a single head. 15 samples were randomly selected from 5 tasks and 3 different context lengths in the RULER benchmark. The visualization shows a portion of the attention heatmap from Layer 3, Head 7 of the LLama3-8B model. For each plot, the first 10 rows correspond to the attention heatmap of the last 10 query tokens during the pre-filling stage, while the remaining rows correspond to the query tokens during the decoding stage.}
%     \vspace{-0.3cm}
%     \label{fig:tasks_lengths}
% \end{center}
% \end{figure}

% \begin{figure}[t]
% \begin{center}
%     \includegraphics[width=0.9\textwidth]{Figures/Similar_focus.pdf}
%     \caption{Similar attention focus across different consistent heads. The attention heatmaps for several heads with consistent patterns in Longchat-7B are presented, applied to the same input. The highlighted regions in these heatmaps show a strong similarity.}
%     \vspace{-0.3cm}
%     \label{fig:similar_focus}
% \end{center}
% \end{figure}

% \begin{figure}[t]
% \begin{center}
%     \includegraphics[rwidth=0.45\textwidth]{Figures/latency_breakdown.pdf}
%     \caption{Latency breakdown of attention layer under different adaptive head ratios. Latencies for consistent heads, adaptive heads, and other operations (index extraction and merging) when tested on LLaMA-3-8B with an input length of 128K. The values represent the total latency in the attention layer for generating 100 tokens.}
%     \vspace{-0.3cm}
%     \label{fig:latency_breakdown}
% \end{center}
% \end{figure}

% \input{Tables/task_aware}






\end{document}
